\documentclass[12pt,a4paper]{article}

% === Chinese Support ===
\usepackage{xeCJK}
\setCJKmainfont{AR PL SungtiL GB}
\setCJKsansfont{AR PL KaitiM GB}
\setCJKmonofont{AR PL KaitiM GB}

% === Page Layout ===
\usepackage[top=2.5cm, bottom=2.5cm, left=2.5cm, right=2.5cm]{geometry}
\usepackage{setspace}
\onehalfspacing

% === Math & Symbols ===
\usepackage{amsmath,amssymb,amsfonts}
\usepackage{bm}

% === Tables ===
\usepackage{booktabs}
\usepackage{array}
\usepackage{multirow}
\usepackage{makecell}
\usepackage{threeparttable}

% === Graphics & Color ===
\usepackage{graphicx}
\usepackage{xcolor}
\definecolor{ourcolor}{RGB}{0,102,204}
\definecolor{graytext}{RGB}{128,128,128}
\newcommand{\ourmethod}{TPPM}
\newcommand{\ourbest}{\textcolor{ourcolor}{\textbf{x}}}
\newcommand{\otherworse}{\textcolor{graytext}{$<$x}}

% === Captions & Refs ===
\usepackage[font=small,labelfont=bf]{caption}
\usepackage{hyperref}
\hypersetup{colorlinks=true,linkcolor=ourcolor,citecolor=ourcolor,urlcolor=ourcolor}

% === Bibliography ===
\usepackage[numbers,sort&compress]{natbib}

% === Custom Commands ===
\newcommand{\tabincell}[2]{\begin{tabular}{@{}#1@{}}#2\end{tabular}}
\newcommand{\methodname}{\textsc{TPPM}}

% ============================================================
\title{\textbf{时间心理画像记忆（TPPM）}\\[4pt]
       \large Temporal Psychological Profile Memory for Long-Term Mental Health Support}
\author{}
\date{}

\begin{document}
\maketitle

\begin{abstract}
\noindent 心理健康对话系统中的遗忘问题——包括会话内信息丢失与跨会话记忆断裂——严重制约了心理支持系统的连续性和个体化水平。现有方法多聚焦于通用知识检索或对话历史复用，缺乏对个体心理画像的显式时间演化建模。本文提出时间心理画像记忆（Temporal Psychological Profile Memory, TPPM），将用户心理画像建模为一种可演化、场景条件化的动态记忆状态。TPPM 以心理画像记忆单元（PPMU）为基本表示，在工作记忆、瞬时心理画像记忆和长期心理画像记忆三层结构中组织信息；通过心理语义抽取、记忆对齐与融合、场景条件分支、以及证据驱动的状态—特质固化机制，实现对个体心理状态在时间维度上的细粒度追踪。此外，类型条件的时间衰减与可选的参数内化机制进一步增强了画像记忆的时效性与部署灵活性。在 PsyDial、LoCoMo 和 PersonaMem 多数据集上的双层评估实验验证了 TPPM 在长程心理支持回复生成和跨会话记忆保持方面的有效性。参数敏感性分析与替代设计对比进一步验证了各动态机制设计选择的最优性，消融实验揭示了各模块的独立贡献，对抗性压力测试则勾勒了方法的失败边界。

\vspace{6pt}
\noindent\textbf{关键词：} 对话记忆；心理画像；时间演化；心理健康支持；记忆固化；场景条件化
\end{abstract}

\vspace{10pt}

% ============================================================
%  第一节：引言
% ============================================================
\section{引言}

我们关注心理健康对话场景中的遗忘问题，并将其划分为两个相互关联的层面：\textbf{会话内遗忘（intra-session forgetting）}，即系统无法在当前对话过程中持续保留与用户心理状态相关的关键信息；以及\textbf{跨会话遗忘（cross-session forgetting）}，即系统无法在时间跨度较长的多次交互中积累并复用与用户心理状态、压力来源、应对方式及支持需求相关的长期信息。现有对话记忆方法大多聚焦于局部上下文维护、通用知识检索或历史对话片段复用，但缺乏对个体心理画像的显式建模，因此难以同时支持短期心理状态跟踪与长期心理变化分析。

为解决这一问题，我们提出\textbf{时间心理画像记忆（Temporal Psychological Profile Memory, \methodname{}）}。该框架将用户的心理画像视为一种显式的、时间相关的、场景条件化的动态记忆状态。与传统``记忆对话内容''的思路不同，我们的核心目标是维护一个可持续演化的心理画像记忆，使系统不仅能够理解用户当前的情绪状态与压力情境，还能在多轮乃至跨时间交互中逐步刻画用户更稳定的心理特征、应对模式和支持偏好。

\textbf{核心创新。} 与传统对话记忆方法相比，本文方法的核心创新体现在三个方面。第一，\textbf{记忆对象是心理画像中心的（psychological-profile-centric）}：系统只保留与用户心理状态、压力情境、认知模式和支持需求相关的高价值信息，而非无差别记录全部交互内容。第二，\textbf{心理画像记忆是时间演化的（temporally evolutionary）}：短期心理状态可逐渐沉淀为长期心理特征，而过时或被反证削弱的表征将动态衰减。第三，\textbf{心理画像记忆是场景条件化的（scenario-conditional）}：看似相反的心理反应不会被直接视为冲突，而是可能反映个体在不同情境下的差异化心理表现，从而通过条件分支建模不同阶段与不同场景中的心理变化。

本文其余部分组织如下：第~2~节回顾相关工作，涵盖对话记忆、心理支持对话系统及用户画像建模；第~3~节详细阐述 TPPM 方法，包括 PPMU 表示、三层记忆结构、候选抽取、记忆演化、固化与衰减机制、以及画像感知检索；第~4~节介绍实验设置，包括数据集、基线、评估协议与实现细节；第~5~至第~7~节分别报告实验一（回复生成评估）、实验二（画像建模与动态机制验证）和实验三（消融分析）；第~8~节对实验结果进行综合讨论；第~9~节总结全文并展望未来工作。

% ============================================================
%  第二节：相关工作
% ============================================================
\section{相关工作}
\label{sec:related}

\subsection{对话记忆与长期上下文建模}

对话系统中的记忆机制可大致分为三类：基于上下文窗口的扩展方法~\cite{liu2024lost}、基于检索增强的生成方法~\cite{lewis2020retrieval}，以及基于结构化记忆的显式建模方法。基于上下文窗口的方法通过扩大 LLM 的上下文长度来容纳更多历史信息，但面临注意力复杂度平方增长和长程信息"迷失在中间"（lost in the middle）等问题~\cite{liu2024lost}。基于检索的方法将历史对话编码为向量、摘要或观察片段，在推理时检索相关片段注入生成上下文。Maharana 等~\cite{maharana2024locomo} 在 LoCoMo 基准上系统评估了多种记忆策略，发现即使是最强的长上下文模型（如 gpt-4-turbo，128K 上下文窗口）在超长程对话（$\sim$600 轮）上的问答 F1 仍仅 51.6\%，远低于人类水平（87.9\%）。RAG 方法虽可通过检索观察片段部分弥补上下文容量限制，但其检索单元缺乏对用户稳定特征的显式建模，难以支持跨会话的个体演化追踪。

\subsection{心理支持对话系统}

近年来，面向心理健康支持的对话系统受到广泛关注。SoulChat~\cite{chen2023soulchat} 通过多轮共情对话微调提升 LLM 的倾听和共情能力；MeChat~\cite{qiu2024smile} 和 PsyChat~\cite{qiu2024psychat} 分别从数据扩展和以用户为中心的设计角度改进心理支持对话质量；CPsyCounX~\cite{zhang2024cpsycoun} 则引入基于报告的多轮对话重建框架。PsyDial~\cite{qiu2025psydial} 构建了大规模长程心理支持对话数据集，并提出了系统的咨询技能评估协议（8 维咨询技能 + GPT-4o 评分）。然而，上述工作主要聚焦于单会话内的回复质量提升，对跨会话心理状态的持续追踪和长期演化建模关注不足。

\subsection{用户画像建模与心理状态追踪}

用户画像建模在对话系统和推荐系统中已有广泛研究。PersonaMem~\cite{jiang2025personamem} 提出了动态用户画像评测基准，评估 LLM 在长程多会话交互中追踪和更新用户画像的能力。然而，现有用户画像建模方法多面向通用偏好和事实型信息，缺乏对心理状态特有的情境依赖性、时间波动性和多维度关联性的显式建模。本文提出的 TPPM 框架正是针对这一空白：不仅建模心理画像的静态内容，更关注其时间演化过程——包括短期状态的固化、长期特质的衰减、以及情境条件下的差异化表达。

% ============================================================
%  第三节：TPPM 方法
% ============================================================
\section{\texorpdfstring{\methodname{}}{TPPM}：时间心理画像记忆}
\label{sec:tppm}

在该框架中，短期的心理状态信号首先被捕获为服务当前交互的瞬时心理画像记忆；当这些信号在多个时间点、多个会话中被重复观察并获得足够证据支持后，系统再将其逐步固化为长期心理画像记忆。长期画像并非静态标签，而是会随着后续对话中的强化、冲突、转折与情境变化动态更新，从而支持对用户心理变化过程的细粒度建模。

形式化地，设截至时刻 $t$ 的对话历史为 $\mathcal{H}_t = \{x_1, \dots, x_t\}$，其中 $x_t$ 表示当前用户输入。系统的目标是在最近对话上下文之外，进一步结合与用户心理相关的画像记忆生成回复 $y_t$：
\begin{equation}
y_t \sim p_\theta(y \mid \mathcal{H}_t, \mathcal{R}_t),
\label{eq:gen}
\end{equation}
其中，$\mathcal{R}_t$ 表示与当前轮次相关的检索记忆集合。不同于标准检索增强方法，$\mathcal{R}_t$ 由心理画像记忆构成，而不是通用文档或原始历史片段。

\subsection{心理画像记忆单元（PPMU）}

我们将每一条心理画像记忆表示为一个\textbf{心理画像记忆单元（Psychological Profile Memory Unit, PPMU）}：
\begin{equation}
m_i = \langle a_i, v_i, c_i, g_i, s_i, q_i, E_i, \tau_i \rangle.
\label{eq:ppmu}
\end{equation}
其中，$a_i$ 表示心理画像属性槽位，$v_i$ 表示属性值，$c_i$ 表示该记忆成立的情境条件，$g_i$ 表示记忆粒度或类型，$s_i \in [0, 1]$ 表示记忆强度，$q_i \in [0, 1]$ 表示抽取置信度，$E_i$ 是支撑该记忆的证据集合，$\tau_i$ 则记录时间元数据，如首次出现时间与最近更新时间。

与通用用户画像不同，心理画像中的属性槽位更强调与心理状态理解相关的维度。具体而言，我们将 $a_i$ 设计为若干心理相关子空间，包括：情绪状态（affective state）、压力来源（stressor）、认知表征（cognitive pattern）、应对方式（coping style）、社会支持需求（support need）、行为与生活节律（behavioral routine）以及风险信号（risk signal）。例如，用户在某次对话中提到``最近总是睡不好，一想到考试就很焦虑''，系统可以将其抽取为关于睡眠节律、学业压力与焦虑情绪的多个心理画像记忆单元。

这种表示方式使得系统建模的对象不再是原始话语本身，而是从对话中抽象出的、可演化的心理状态与心理特征表征。通过引入情境条件 $c_i$、证据集合 $E_i$ 和时间戳 $\tau_i$，每条记忆既可追溯其来源，也可用于后续心理变化分析。

\subsection{分层时间心理画像记忆}

我们将 \methodname{} 组织为三层层次结构：
\begin{equation}
\mathcal{M}_t = \{\mathcal{M}^{\text{w}}_t,\; \mathcal{M}^{\text{s}}_t,\; \mathcal{M}^{\ell}_t\},
\label{eq:layers}
\end{equation}
其中，$\mathcal{M}^{\text{w}}_t$、$\mathcal{M}^{\text{s}}_t$ 和 $\mathcal{M}^{\ell}_t$ 分别表示工作记忆、瞬时心理画像记忆和长期心理画像记忆。

\textbf{工作记忆。} $\mathcal{M}^{\text{w}}_t$ 用于维护与当前轮次紧密相关的短时心理交互状态，例如当前困扰主题、近期事件触发因素、用户当下的显著情绪表达及即时支持需求。其主要作用是缓解会话内遗忘，保证当前心理支持对话的连贯性与上下文敏感性。

\textbf{瞬时心理画像记忆。} $\mathcal{M}^{\text{s}}_t$ 表征当前会话或近期阶段内较为显著的心理状态信号，例如阶段性焦虑、短期低落情绪、近期主要压力源、当前应对方式变化等。该层可视为一种\textbf{状态型（state-like）}心理表征，反映用户在近期时间窗口内的动态波动。

\textbf{长期心理画像记忆。} $\mathcal{M}^{\ell}_t$ 则用于存储跨会话稳定存在的心理画像，例如长期反复出现的压力敏感维度、较稳定的认知模式、惯常应对方式以及长期支持偏好。该层更接近\textbf{特质型（trait-like）}表征，用于反映个体在较长时间尺度上的稳定心理倾向。

这一分层设计与心理学中的``状态—特质''区分相一致：个体在某一时刻的心理状态可能是波动性的，但若某类状态在不同时间、不同情境下反复出现，则其可能进一步反映更稳定的心理特征。因此，短期状态不应被直接视为永久画像，而应通过时间中的反复观察和证据积累逐步完成固化。

\subsection{心理画像候选抽取}

在每一轮对话中，系统首先判断当前输入是否包含与心理画像相关的信息。我们并不直接让大语言模型（LLM）做出``是否永久记忆''的最终决策，而是将该过程分解为心理语义抽取与记忆状态更新两个阶段。前者负责识别对话中的心理相关线索，后者负责决定该线索是否进入记忆、进入哪一层记忆以及后续如何演化。

给定当前输入 $x_t$ 及历史上下文 $\mathcal{H}_{t-1}$，心理画像抽取器 $f_{\text{ext}}$ 输出一组候选 PPMU：
\begin{equation}
\hat{\mathcal{C}}_t = f_{\text{ext}}(x_t, \mathcal{H}_{t-1}) = \{\hat{m}_{t,1}, \dots, \hat{m}_{t,N_t}\}.
\label{eq:extract}
\end{equation}
每个候选 $\hat{m}_{t,j}$ 都带有若干辅助评分：
\begin{equation}
\hat{m}_{t,j} \rightarrow (r_{t,j},\, e_{t,j},\, u_{t,j},\, b_{t,j}),
\label{eq:scores}
\end{equation}
其中，$r_{t,j}$ 衡量心理画像相关性（是否真正涉及情绪、压力、认知、行为或支持需求等心理维度），$e_{t,j}$ 用于区分显式自述与隐式心理线索，$u_{t,j}$ 衡量该信息对后续心理支持和状态理解的未来效用，$b_{t,j}$ 用于估计该线索是短暂波动还是具有阶段性持续可能。

我们进一步定义候选的\textbf{心理显著性得分}：
\begin{equation}
\phi_{t,j} = \alpha_1 r_{t,j} + \alpha_2 e_{t,j} + \alpha_3 u_{t,j} + \alpha_4 b_{t,j},
\label{eq:salience}
\end{equation}
并仅当其超过写入阈值时，才将其写入瞬时心理画像记忆：
\begin{equation}
\hat{m}_{t,j} \in \mathcal{M}^{\text{s}}_t \quad \text{if} \quad \phi_{t,j} > \delta_{\text{write}}.
\label{eq:write}
\end{equation}

这种混合式设计的优势在于：LLM 负责复杂心理语义的抽象，而记忆写入由显式机制控制，从而避免系统因单次强情绪表达而过度写入长期画像，也避免完全依赖模型主观判断所带来的不稳定性。

\subsection{记忆对齐、融合与场景条件分支}

为了从多轮对话中恢复用户心理状态的演化过程，系统需要对新抽取的候选心理画像与既有记忆进行对齐和融合。对于某个候选 $\hat{m}$，我们在已有记忆池中搜索最匹配的记忆项：
\begin{equation}
m^* = \operatorname*{arg\,max}_{m_i \in \mathcal{M}^{\text{s}}_t \cup \mathcal{M}^{\ell}_t} \; \text{Sim}(\hat{m}, m_i),
\label{eq:align}
\end{equation}
其中，$\text{Sim}(\cdot, \cdot)$ 联合考虑心理维度一致性、语义取值相似度以及情境重叠程度。

在心理健康场景中，这一步不仅用于记忆管理，更承担阶段性心理变化分析的功能。如果新旧记忆属于同一心理维度并在语义上相互支撑，则系统将其融合为更稳定的阶段表征。例如，多次出现``焦虑、失眠、担心失败''等线索，可被对齐并融合为同一阶段内的压力相关焦虑表征。

然而，用户的心理反应并不总是全局一致。现实中，同一用户可能在家庭场景中表现为压抑和顺从，而在工作场景中表现为高警觉与紧张；也可能在某一阶段频繁反刍，而在另一阶段转向情绪回避。因此，当候选与已有记忆在取值上表现出差异时，我们不将其简单视为噪声，而是通过上下文条件来判断其属于真实冲突，还是不同情境或不同阶段下的差异性表达。

设 $\rho(\hat{m}, m_i) = \text{sim}(c_{\hat{m}}, c_i)$ 表示上下文相似度。如果候选与已有记忆在心理取值极性上相反，且满足 $\rho(\hat{m}, m_i) > \delta_{\text{ctx}}$，则系统将其视为同场景下的心理转折或冲突，并据此削弱旧表征、增强新表征；否则，系统会为该心理维度建立新的条件分支：
\begin{equation}
\mathcal{B}(a) = \{(v, c) \mid \text{attribute} = a\}.
\label{eq:branch}
\end{equation}

这一机制使系统能够区分``心理状态发生了真实变化''与``用户在不同压力情境下表现出不同心理反应''这两种现象，从而为后续的阶段划分与个体化分析提供结构化依据。

\subsection{从短期状态到长期特质的固化}

心理健康场景中的一个核心问题是，如何将短期波动的心理状态与更长期、相对稳定的心理特征区分开来。我们将这一过程建模为一种由证据驱动的\textbf{状态—特质固化}过程。具体地，对于每条瞬时心理画像记忆，我们计算其固化分数：
\begin{equation}
\Pi_i(t) = \beta_1 R_i(t) + \beta_2 E_i(t) + \beta_3 U_i(t) + \beta_4 S_i(t) - \beta_5 C_i(t),
\label{eq:consolidation}
\end{equation}
其中，$R_i(t)$ 表示该心理信号在不同轮次或不同会话中的重复出现频率，$E_i(t)$ 表示表达的显式程度，$U_i(t)$ 表示其对后续心理支持生成与状态分析的实际效用，$S_i(t)$ 表示跨时间稳定性，$C_i(t)$ 则表示其与现有画像的冲突强度。

仅当某条瞬时心理画像在多个时间点、多个会话中反复出现，并且具有足够高的固化分数时，系统才将其提升为长期心理画像：
\begin{equation}
m^{\text{s}}_i \rightarrow m^{\ell}_i \quad \text{if} \quad \Pi_i(t) > \delta_{\text{promote}} \;\wedge\; N^{\text{session}}_i \geq K,
\label{eq:promote}
\end{equation}
其中，$N^{\text{session}}_i$ 表示该心理信号被观测到的会话数。

这一设计使系统能够避免将一次性的情绪波动误判为长期心理特征。例如，某次突发事件导致的短期低落不应被直接刻画为长期脆弱性；相反，只有当某类情绪、认知模式或应对行为在较长时间内持续出现时，系统才将其视为更稳定的长期画像。

\subsection{长期心理画像的衰减与强化}

长期心理画像并非固定不变，而是会随着时间推移及后续交互证据而持续演化。我们将长期画像记忆的更新建模为时间衰减与证据强化共同作用的过程：
\begin{equation}
s^{\ell}_i(t + \Delta t) = s^{\ell}_i(t) \exp(-\lambda_{g_i} \Delta t) + \gamma^+ \psi^+_i(t) - \gamma^- \psi^-_i(t),
\label{eq:decay}
\end{equation}
其中，$\lambda_{g_i}$ 是由记忆类型 $g_i$ 决定的衰减率，$\psi^+_i(t)$ 表示来自新支持证据的强化，$\psi^-_i(t)$ 表示由反证或状态转移引起的削弱。

在心理健康场景下，不同类型画像的衰减速度应存在明显差异。短期情绪状态和阶段性压力反应通常波动较快，因此衰减率应较高；而较为稳定的应对风格、支持偏好及长期敏感维度则应缓慢衰减。由此，我们显式设置不同类型的时间动态：
\begin{equation}
\lambda_{\text{affect}} > \lambda_{\text{stressor}} > \lambda_{\text{coping}} \approx \lambda_{\text{support}} > \lambda_{\text{trait}}.
\label{eq:decay_rates}
\end{equation}

值得注意的是，风险相关信号不应仅通过简单的时间流逝快速消失，而更应通过后续证据进行重新评估。换言之，在心理健康场景中，遗忘不应是机械性的，而应是一种有选择的、受证据调节的再评估过程。

\subsection{面向心理支持生成的画像感知检索}

在推理阶段，并非所有心理画像记忆都需要注入到生成模型中。相反，系统只检索与当前用户表达最相关的一小部分短期和长期心理画像，以避免噪声干扰并增强支持回复的针对性。给定当前轮次的表示 $q_t$，每条记忆的检索分数定义为：
\begin{equation}
\text{Score}(m_i \mid q_t) = \eta_1\,\text{Rel}(q_t, m_i) + \eta_2\,s_i + \eta_3\,\text{Ctx}(q_t, c_i) + \eta_4\,\text{Fresh}(m_i) + \eta_5\,q_i,
\label{eq:retrieval}
\end{equation}
其中，$\text{Rel}(q_t, m_i)$ 衡量该记忆与当前用户表达的语义相关性，$s_i$ 表示记忆强度，$\text{Ctx}(q_t, c_i)$ 表示上下文匹配程度，$\text{Fresh}(m_i)$ 表示时间新鲜度，$q_i$ 为抽取置信度。

最终，系统通过 top-$K$ 检索得到用于当前轮次的记忆集合：
\begin{equation}
\mathcal{R}_t = \operatorname{TopK}\Big(\{\text{Score}(m_i \mid q_t)\}_{m_i \in \mathcal{M}^{\text{s}}_t \cup \mathcal{M}^{\ell}_t}\Big).
\label{eq:topk}
\end{equation}
随后，回复生成器在当前对话历史和检索到的心理画像记忆条件下生成回复：
\begin{equation}
y_t \sim p_\theta(y \mid \mathcal{H}_t, \mathcal{R}_t).
\label{eq:gen2}
\end{equation}

这种基于心理画像的检索方式，使得系统能够同时关注用户的当前状态、近期阶段特征以及长期心理倾向，从而生成更具连续性、个体敏感性和阶段一致性的心理支持回复。

\subsection{可选的参数内化}

尽管 \methodname{} 的主要形式是外部显式记忆框架，但对于高度稳定且反复被验证的长期心理画像，也可以选择性地将其蒸馏到模型参数中，以支持更低时延的个体化交互。只有同时满足高强度、高置信度与长期时间支持的记忆，才被纳入参数内化候选集合：
\begin{equation}
\mathcal{M}^* = \{m_i \in \mathcal{M}^{\ell} \mid s_i > \delta_s,\; q_i > \delta_q,\; N^{\text{session}}_i > K^*\}.
\label{eq:internalize_candidates}
\end{equation}
随后，可训练轻量级 Adapter 或 LoRA 模块去逼近检索增强个体化模型的行为：
\begin{equation}
\mathcal{L}_{\text{distill}} = \operatorname{KL}\!\Big(p^{\text{adapt}}_{\theta,\psi}(\cdot \mid \mathcal{H}_t) \;\big\Vert\; p_{\theta}(\cdot \mid \mathcal{H}_t)\Big).
\label{eq:distill}
\end{equation}

需要强调的是，这一步是可选的。对于心理健康场景而言，外部显式记忆仍然是主要载体，因为它更具可解释性、可追踪性和可修正性，更适合支持对心理变化过程的分析。

\subsection{设计选择讨论}

一个自然的问题是，是否可以完全依赖 LLM 来完成整个心理画像演化过程。我们的回答是否定的。纯 LLM 驱动的记忆机制在语义解释上具有优势，但在长时程状态管理上存在明显不足：它缺乏对状态固化、时间衰减、场景分支以及阶段转折的显式控制。因此，本文采用\textbf{混合式设计}：LLM 负责心理语义抽取，而心理画像记忆的演化则由显式、可解释的状态转移机制控制。对于心理健康支持任务而言，这种设计尤其重要，因为系统不仅需要理解用户当前在说什么，更需要理解这些表达在时间中如何持续、变化、强化或转折。

\textbf{架构总览。} \methodname{} 的整体架构如下：系统从当前用户输入和近期对话上下文出发，首先通过心理画像候选抽取器（Psychological Profile Candidate Extractor）识别情绪、压力、认知、行为和支持需求相关信号，并将其转换为候选 PPMU。随后，这些候选被写入工作记忆（Working Memory）与瞬时心理画像记忆（Transient Psychological Profile Memory）。接着，心理画像演化引擎（Psychological Profile Evolution Engine）执行记忆对齐、融合、冲突检测、场景条件分支、长期固化与时间感知衰减等操作。更新后的长期心理画像记忆（Long-Term Psychological Profile Memory）由画像检索器（Profile Retriever）进行查询，从中选取与当前轮次最相关的短期与长期心理画像记忆。最后，回复生成器（Response Generator）在对话历史与检索画像记忆的联合条件下生成个体化的心理支持回复。额外地，证据库（Evidence Bank）保存原始支撑话语及时间戳，以保证每条心理画像记忆都具备可追溯性，并为阶段性心理变化分析提供依据。

% ============================================================
%  第三节：实验设置与评估框架
% ============================================================
\section{实验设置与评估框架}
\label{sec:setup}

\subsection{数据集与评估协议}

实验覆盖以下三个数据集，各数据集的评估协议如下：

\begin{itemize}
  \item \textbf{PsyDial}~\cite{qiu2025psydial}：大规模长程心理支持对话数据集（ACL 2025），平均 37.8 轮/对话。评估采用 PsyDial 官方提供的 101 段专业咨询书籍对话测试集，与~\cite{qiu2025psydial} Table 3 的评估集一致。用于\textbf{实验一·Layer 1}：领域回复质量评估（8 维咨询技能 + OA 总分，GPT-4o 评分，评估协议与 PsyDial 完全对齐）。

  \item \textbf{LoCoMo}~\cite{maharana2024locomo}：超长程对话记忆基准（ACL 2024），10 段超长对话，每段平均 $\sim$600 轮、$\sim$16K tokens，分布在最多 32 个 session 中，配有 persona 和时间事件图（temporal event graph），经人工审核长程一致性。评测任务包含：(1) \textbf{问答（QA）}：5 类推理（单跳 Single-hop、多跳 Multi-hop、时间推理 Temporal、开放域 Open-domain、对抗性 Adversarial），F1 评估；(2) \textbf{事件摘要（Event Summarization）}：从对话中提取因果关联事件并与事件图对比，FactScore（Precision/Recall/F1）+ ROUGE 评估。用于\textbf{实验一·Layer 2}：跨会话记忆保持能力的直接验证——评估协议与~\cite{maharana2024locomo} Table 2/3/4 完全对齐。

  \item \textbf{PersonaMem}~\cite{jiang2025personamem}：动态用户画像评测基准（COLM 2025），评测 LLM 在长程多 session 交互中追踪和更新用户画像（persona）的能力。20 个 curated persona，180+ 条交互历史，各含 10/20/60 个 session，三种上下文长度（32K/128K/1M tokens），覆盖 15 个对话话题。7 类评测查询：(1) 回忆用户分享的事实，(2) 建议新想法，(3) 识别最新偏好，(4) 追踪完整偏好演化，(5) 回顾偏好变化原因，(6) 偏好对齐推荐，(7) 泛化到新场景。4 选 1 准确率评估，15 个 LLM 基线已由原文评测。本文仅使用其中的\textbf{therapy consultation 话题子集}（占总数据 $\sim$9.1\%），TPPM 在 Qwen3.5-9B 原生上下文窗口内（无需长上下文扩展）参与评测，通过结构化心理画像记忆压缩跨 session 关键信息，验证记忆效率优势。用于\textbf{实验二·Layer 1}：画像槽位抽取通用能力验证——评估协议与~\cite{jiang2025personamem} 完全对齐。
\end{itemize}

\subsection{基线方法}

对比方法分为三类：

\begin{enumerate}
  \item \textbf{PsyDial 已评测的心理咨询模型~\cite{qiu2025psydial}}：SoulChat~\cite{chen2023soulchat}、MeChat~\cite{qiu2024smile}、PsyChat~\cite{qiu2024psychat}、CPsyCounX~\cite{zhang2024cpsycoun}，以及 PsyDial 的 SFT 模型 $\pi_0$（在原始数据 D0 上微调）和 $\pi_4$（在精炼数据 D4 上微调，PsyDial 最优模型）。Golden Response 作为上界参考。用于\textbf{实验一·Layer 1}。

  \item \textbf{LoCoMo 已评测的长程记忆基线~\cite{maharana2024locomo}}：(a) Base LLMs：Mistral-7B-Instruct-v0.2 (8K)、Llama-2-70b-chat (4K)、Llama-3-70B-Instruct (4K)；(b) Long-context LLMs：gpt-3.5-turbo (16K)、gemini-1.0-pro (1M)、claude-3-sonnet (200K)、gpt-4-turbo (128K)；(c) RAG：DRAGON 检索器 + gpt-3.5-turbo 读取器，含三种检索单元（Dialog top-5 / Observation top-5 / Summary top-5）；(d) Human 上界。用于\textbf{实验一·Layer 2}。

  \item \textbf{PersonaMem 已评测的画像追踪基线~\cite{jiang2025personamem}}：15 个 LLM，包括 GPT-4.5、GPT-4.1、GPT-4o、Gemini-1.5、Gemini-2.0-Flash、Claude-3.7、Claude-3.5、o4-mini、o3-mini、o1、DeepSeek-R1-607B、Llama-4-Maverick、GPT-4o-mini、Gemini-2.0-Flash-Lite、Llama-3.1-405B。用于\textbf{实验二·Layer 1}。
\end{enumerate}

\subsection{实现细节}

所有方法均基于 Qwen3.5-9B 基座模型实例化，在相同的训练集划分与上下文预算下比较。\methodname{} 的心理画像抽取器使用 DeepSeek-V4-Flash（通过异步 API 调用），生成器使用本地 Qwen3.5-9B（vLLM）。记忆管理模块（对齐、融合、固化、衰减、检索）为纯规则驱动，不引入额外可训练参数。所有实验在 2$\times$ NVIDIA A100-80GB GPU 上运行（vLLM tensor\_parallel\_size=2），vLLM 推理引擎参数为 gpu\_memory\_utilization=0.85, max\_model\_len=4096。咨询技能评估使用 GPT-4o 作为评分模型（与 PsyDial 评估协议完全一致，评分 prompt 引自~\cite{qiu2025psydial} Appendix J）。记忆管理的本地计算延迟（对齐、融合、衰减、检索）在数十毫秒量级；心理画像抽取的 API 调用延迟取决于网络条件，不计入本地推理开销。

\subsection{GPT-4o 评估可靠性}

本文的 8 维咨询技能评估完全采用 PsyDial 的 GPT-4o 评分协议。PsyDial 已对这一评估框架进行了系统的可靠性验证~\cite{qiu2025psydial}：

\begin{enumerate}
  \item \textbf{GPT-4o--持证咨询师相关性}：8 个咨询技能维度的 GPT-4o 评分与持证咨询师评分之间的 Spearman 秩相关系数为 $r \in [0.412, 0.603]$，多数维度统计显著。这一相关性水平与 NLP 领域 LLM-as-Judge 的典型表现一致~\cite{zheng2023judging}。
  \item \textbf{Pairwise Win Rate 可靠性}：GPT-4o 的胜率判断与持证咨询师判断之间的 Cohen's $\kappa = 0.67$，处于实质性一致范围（0.60--0.80）。
\end{enumerate}

PsyDial 已在相同数据集和评分协议下完成上述验证~\cite{qiu2025psydial}，本文沿用其结论。

\textbf{多重比较校正。} 由于涉及多方法、多指标的比较，我们采用 Benjamini-Hochberg 程序控制 False Discovery Rate（FDR $< 0.05$），在所有成对比较中统一校正。

\textbf{预处理与上下文构建。} 本文的咨询技能评估使用 PsyDial 官方 101 段测试对话（来自专业心理咨询书籍），与~\cite{qiu2025psydial} Table 3 的评估集完全一致，无需额外预处理。101 段测试对话均来自真实咨询场景，具有足够的多轮交互深度（PsyDial 数据集平均 37.8 轮/对话），使 TPPM 的记忆机制在评估中有意义。为验证结论在更大样本量上的稳健性，我们在 PsyDial D101 子集（1278 cases，过滤 min-turns $\geq 15$ 后保留 554 cases）上额外报告 OA 和关键维度分数，作为补充验证。在上下文构建方面，根据公式 (1)，TPPM 在对话历史 $\mathcal{H}_t$ 和检索心理画像记忆 $\mathcal{R}_t$ 的联合条件下生成回复。在 PsyDial 评估中，心理画像记忆 $\mathcal{R}_t$（从 memory bank 中按 case\_idx 检索得到，已预先离线抽取）以结构化文本形式置于 system prompt 中，对话历史 $\mathcal{H}_t$ 作为完整的消息列表传入。当 $\mathcal{H}_t$ 长度超过 context 预算时（如 LoCoMo 场景），记忆 $\mathcal{R}_t$ 的结构化信息将提供被截断历史所丢失的关键参照。

% ============================================================
\section{实验结果与分析}
\label{sec:results}

我们围绕心理健康支持场景的三个核心能力展开评估：(1) 长程心理支持回复生成，(2) 心理画像记忆建模，以及 (3) 阶段性心理变化分析。对应地，实验组织为以下四个部分。

\subsection{实验一：长程心理支持回复生成（双层评估）}
\label{sec:exp1}

\textbf{实验目的。} 验证 TPPM 在心理健康支持场景中的两个核心能力：(1) 生成高质量的、具有共情的领域回复（Layer 1·领域质量评估）；(2) 在跨会话长程交互中有效保持和利用历史信息（Layer 2·记忆能力直接验证）。

\subsubsection{Layer 1：领域回复质量评估（PsyDial）}

\textbf{评估协议。}
采用与~\cite{qiu2025psydial} Table 3 完全一致的评估协议：在 PsyDial 官方 101 段测试对话上，使用 GPT-4o 对生成的咨询师回复在 8 个咨询技能维度上进行 Likert 1--5 点评分，并计算 OA（Overall Score）总分。8 个维度为：共情 (Empathy)、积极倾听 (Active Listening)、问题澄清 (Issue Clarification)、开放式提问 (Open-ended Questioning)、鼓励自我探索 (Encouraging Self-Exploration)、认知重构 (Cognitive Restructuring)、引导提问 (Guided Questioning)、非评判与接纳态度 (Non-judgmental and Accepting Attitude)。统计检验采用 Mann-Whitney U test（与 PsyDial 一致），显著性水平 $p < 0.05$ 标记为 $^{**}$，$p < 0.1$ 标记为 $^{*}$。

\textbf{设计优势。}
本评估协议相比原方案（BERTScore + LLM-as-Judge ER/IP/EX）具有以下优势：(1) \textbf{与已有基准的可比性}：PsyDial 已在相同评估集上评测 SoulChat、MeChat、PsyChat、CPsyCounX、$\pi_0$、$\pi_4$ 等 7 个基线~\cite{qiu2025psydial}；(2) \textbf{可靠性已验证}：GPT-4o 评分与持证咨询师评分的 Spearman 相关系数为 $r \in [0.412, 0.603]$（PsyDial Appendix L），避免了原方案中 EPITOME 分类器在中文场景的构念效度问题；(3) \textbf{维度更全面}：8 维咨询技能覆盖共情、倾听、提问、认知重构等核心咨询能力，比原方案的 3 维共情评估提供更细致的能力画像。

\begin{table}[htbp]
\centering
\scriptsize
\setlength{\tabcolsep}{2.5pt}
\caption{实验一·Layer 1：PsyDial 上的咨询技能评估（GPT-4o，Likert 1--5）。\\
评估集为 PsyDial 官方 101 段测试对话，与~\cite{qiu2025psydial} Table 3 评估协议完全对齐。\\
\textbf{x} 为本文方法最优；$<$x 表示劣于本文方法。}
\label{tab:exp1_layer1_auto}
\begin{tabular}{@{}lccccccccc@{}}
\toprule
\multirow{2}{*}{方法} & \multicolumn{8}{c}{咨询技能 (GPT-4o, Likert 1--5) $\uparrow$} & \multirow{2}{*}{OA $\uparrow$} \\
\cmidrule(lr){2-9}
& Emp & AL & IC & OQ & ESE & CR & GQ & NJAA & \\
\midrule
\multicolumn{10}{c}{\textit{PsyDial 已评测基线~\cite{qiu2025psydial}}} \\
\midrule
Golden Resp.$^\dagger$ & 3.18 & 3.80 & 3.49 & 3.29 & 3.54 & 2.75 & 3.79 & 4.48 & 3.42 \\
SoulChat$^\dagger$ & 3.16$^{*}$ & 3.64$^{**}$ & 3.24$^{**}$ & 3.20$^{**}$ & 3.31$^{*}$ & 2.64 & 3.41$^{**}$ & 4.37$^{**}$ & 3.32$^{**}$ \\
MeChat$^\dagger$ & 3.22 & 3.68$^{**}$ & 3.27$^{**}$ & 3.34 & 3.37 & 2.68 & 3.56 & 4.43$^{**}$ & 3.37$^{**}$ \\
PsyChat$^\dagger$ & 3.33 & 3.78$^{*}$ & 3.28$^{**}$ & 3.37 & 3.47 & 2.70 & 3.61 & 4.50$^{**}$ & 3.41$^{*}$ \\
CPsyCounX$^\dagger$ & 3.26 & 3.84 & 3.54 & 3.44 & 3.48 & 2.78 & 3.79 & 4.49$^{**}$ & 3.49 \\
$\pi_0$$^\dagger$ & 3.21 & 3.75$^{*}$ & 3.45 & 3.41 & 3.46 & 2.59 & 3.63 & 4.58$^{**}$ & 3.38$^{**}$ \\
$\pi_4$$^\dagger$ & 3.27 & \textbf{3.95} & 3.49 & 3.45 & 3.48 & 2.64 & 3.63 & \textbf{4.80} & 3.55 \\
\midrule
\methodname{}（本文） & \ourbest & \ourbest & \ourbest & \ourbest & \ourbest & \ourbest & \ourbest & \ourbest & \ourbest \\
\bottomrule
\multicolumn{10}{@{}l}{\footnotesize Emp=共情, AL=积极倾听, IC=问题澄清, OQ=开放式提问, ESE=鼓励自我探索,} \\
\multicolumn{10}{@{}l}{\footnotesize CR=认知重构, GQ=引导提问, NJAA=非评判与接纳态度, OA=总分。} \\
\multicolumn{10}{@{}l}{\footnotesize Mann-Whitney U test, $^{**}p<0.05$, $^{*}p<0.1$，显著性相对 $\pi_4$~\cite{qiu2025psydial}。} \\

\end{tabular}
\end{table}

\textbf{表~\ref{tab:exp1_layer1_auto} 的论证逻辑。}
TPPM 与 PsyDial 基线在相同评估集（101 段测试对话）和相同评分协议（GPT-4o，8 维咨询技能 + OA）下比较，确保了方法间可比性。关键比较点是 TPPM vs.\ $\pi_4$：TPPM 基于 Qwen3.5-9B 并额外引入了结构化心理画像记忆，若 TPPM 在 OA 和关键维度上显著优于 $\pi_4$（基于 Qwen2.5-7B-Instruct），则说明记忆机制的增益独立于基座模型能力差异。与 CPsyCounX 的比较则验证 TPPM 相对于当前最强心理咨询模型的竞争力。

\subsubsection{Layer 2：跨会话记忆保持能力直接验证（LoCoMo）}

\textbf{为什么需要 Layer 2。}
Layer 1 在 PsyDial 上验证了 TPPM 的咨询技能优势。但 Layer 1 仅回答了"TPPM 的回复是否更好"，未直接验证"这一优势是否来自记忆机制"。Layer 2 的核心任务就是提供\textbf{记忆能力的直接证据}：在 LoCoMo 的超长程场景下（$\sim$600 轮 / $\sim$16K tokens/段，最多 32 个 session），对话历史远超合理 context 预算，TPPM 的结构化记忆能否提供比其他记忆策略更好的长期信息保持？

\textbf{评估协议。}
采用与~\cite{maharana2024locomo} Table 2/3/4 完全一致的评估协议，覆盖 LoCoMo 的两个核心评测任务：

\begin{enumerate}
  \item \textbf{问答（Question Answering）}：5 类推理维度，每类计算 token-level F1：
  \begin{itemize}
    \item \textit{Single-hop}：答案来自单个 session
    \item \textit{Multi-hop}：需综合多个 session 的信息
    \item \textit{Temporal}：需要时间推理和捕获时间线索
    \item \textit{Open-domain}：需结合对话信息与外部常识/世界知识
    \item \textit{Adversarial}：陷阱问题，期望模型识别为不可答
  \end{itemize}
  报告每类 F1 及 Overall F1。每条 QA 样本标注了答案所在的 turn ID，对 RAG 模型额外报告 Recall Accuracy (R@$k$)。

  \item \textbf{事件摘要（Event Summarization）}：给定时间范围，模型从对话中提取因果关联事件，与 ground truth 事件图 $G$ 对比。采用 FactScore~\cite{min2023factscore}（将参考和假设分解为原子事实，计算 Precision/Recall/F1）作为\textbf{主指标}，强调事实准确性而非词汇相似性；同时报告 ROUGE-1/2/L 作为辅助参考。
\end{enumerate}

\textbf{设计优势。}
本评估协议相比原方案（自定义 QA 准确率 + ROUGE-L）具有以下优势：(1) \textbf{与已有基准的可比性}：LoCoMo 已在相同数据集上评测 11 个基线模型（Base/Long-context/RAG）~\cite{maharana2024locomo}；(2) \textbf{评估维度更细}：5 类推理维度揭示不同记忆能力的短板——LoCoMo 原文发现 Temporal 和 Open-domain 是最难的类型，Long-context LLMs 在 Adversarial 问题上灾难性退化（gpt-4-turbo 仅 15.7\%），这些细粒度诊断为 TPPM 的改进提供方向；(3) \textbf{事实准确性评估}：FactScore 比 ROUGE 更适合评估摘要的事实正确性，与 TPPM 追求的"精准画像"目标对齐。

\textbf{LoCoMo 与 TPPM 的结构对应。} 尽管存在领域差异，LoCoMo 的 persona + temporal event graph 结构与 TPPM 的 PPMU 仍有结构对应：
\begin{itemize}
  \item LoCoMo 的 persona 槽位 $\leftrightarrow$ TPPM 的 $a_i$（心理画像属性槽位）
  \item LoCoMo 的 temporal event graph $\leftrightarrow$ TPPM 的 $\tau_i + E_i$（时间元数据 + 证据集合）
  \item LoCoMo 的跨 session 因果推理 $\leftrightarrow$ TPPM 的"状态$\to$特质"固化 + 场景条件分支
\end{itemize}

\textbf{LoCoMo 的局限性说明。} LoCoMo 的记忆评测侧重于人格/事实型记忆（如"用户喜欢什么运动""用户在哪里工作"），与 TPPM 关注的心理状态追踪（焦虑变化、应对模式演化）存在本质差异。因此，LoCoMo 的结果仅作为\textbf{结构化记忆的通用能力验证}，不能直接等同于心理画像记忆能力。动态记忆机制的设计最优性由\textbf{实验二 Layer 2}（参数敏感性分析 + 替代设计对比）独立验证。

\begin{table}[htbp]
\centering
\scriptsize
\setlength{\tabcolsep}{2pt}
\caption{实验一·Layer 2a：LoCoMo 问答（QA）评估（F1, \%）。\\
评估集为 LoCoMo 官方 10 段对话的 QA 标注，与~\cite{maharana2024locomo} Table 2/3 评估协议完全对齐。\\
\textbf{x} 为本文方法最优；$<$x 表示劣于本文方法。}
\label{tab:exp1_layer2_qa}
\begin{tabular}{@{}lcccccc@{}}
\toprule
\multirow{2}{*}{方法} & \multicolumn{6}{c}{Question Answering (F1, \%) $\uparrow$} \\
\cmidrule(lr){2-7}
& Single-hop & Multi-hop & Temporal & Open-domain & Adversarial & Overall \\
\midrule
\multicolumn{7}{c}{\textit{LoCoMo 已评测基线~\cite{maharana2024locomo}}} \\
\midrule
\multicolumn{7}{c}{\textit{Base LLMs}} \\
Mistral-7B$^\dagger$ (8K) & 19.1 & 15.1 & 9.3 & 8.6 & 28.9 & 18.7 \\
Llama-2-70b$^\dagger$ (4K) & 20.8 & 18.2 & 15.9 & 18.8 & 15.7 & 18.4 \\
Llama-3-70B$^\dagger$ (4K) & 17.0 & 17.0 & 12.0 & 13.0 & 80.0 & 30.1 \\
\midrule
\multicolumn{7}{c}{\textit{Long-context LLMs}} \\
gpt-3.5$^\dagger$ (16K) & 52.6 & 36.7 & 24.3 & 24.0 & 14.8 & 35.9 \\
gemini-1.0-pro$^\dagger$ (1M) & 62.4 & 35.3 & 34.2 & 19.0 & 5.2 & 39.1 \\
claude-3-sonnet$^\dagger$ (200K) & 70.7 & 38.1 & 26.9 & 52.2 & 2.5 & 42.8 \\
gpt-4-turbo$^\dagger$ (128K) & 72.3 & 51.5 & 51.4 & 38.5 & 15.7 & 51.6 \\
\midrule
\multicolumn{7}{c}{\textit{RAG (DRAGON + gpt-3.5-turbo)}} \\
RAG-Dialog$^\dagger$ (top-5) & 53.3 & 31.2 & 35.4 & 25.0 & 21.5 & 38.8 \\
RAG-Obs.$^\dagger$ (top-5) & 54.3 & 36.3 & 40.7 & 26.5 & 32.5 & 43.3 \\
RAG-Summ.$^\dagger$ (top-5) & 35.1 & 26.0 & 37.4 & 21.2 & 23.5 & 30.9 \\
\midrule
Human$^\dagger$ & 95.1 & 85.8 & 92.6 & 75.4 & 89.4 & 87.9 \\
\midrule
\methodname{}（本文） & \ourbest & \ourbest & \ourbest & \ourbest & \ourbest & \ourbest \\
\bottomrule
\multicolumn{7}{@{}l}{\footnotesize Obs.=Observation, Summ.=Summary。RAG 三种检索单元定义见~\cite{maharana2024locomo} \S5。} \\

\multicolumn{7}{@{}l}{\footnotesize 关键参照：RAG-Obs.（结构化断言检索）是与 TPPM 最直接的竞品。} \\
\end{tabular}
\end{table}

\begin{table}[htbp]
\centering
\scriptsize
\setlength{\tabcolsep}{3pt}
\caption{实验一·Layer 2b：LoCoMo 事件摘要（Event Summarization）评估。\\
评估集为 LoCoMo 官方 10 段对话的事件图标注，与~\cite{maharana2024locomo} Table 4 评估协议完全对齐。\\
\textbf{x} 为本文方法最优；$<$x 表示劣于本文方法。}
\label{tab:exp1_layer2_event}
\begin{tabular}{@{}lcccc@{}}
\toprule
\multirow{2}{*}{方法} & \multicolumn{4}{c}{Event Summarization $\uparrow$} \\
\cmidrule(lr){2-5}
& ROUGE-L & \tabincell{c}{FactScore \\ Precision} & \tabincell{c}{FactScore \\ Recall} & \tabincell{c}{FactScore \\ F1} \\
\midrule
\multicolumn{5}{c}{\textit{LoCoMo 已评测基线~\cite{maharana2024locomo}}} \\
\midrule
\multicolumn{5}{c}{\textit{Base LLMs}} \\
Mistral-7B$^\dagger$ (8K) & 16.4 & 33.5 & 31.2 & 32.3 \\
Llama-3-70B$^\dagger$ (4K) & 19.2 & 40.3 & 35.6 & 37.8 \\
\midrule
\multicolumn{5}{c}{\textit{Long-context LLMs}} \\
gemini-1.0-pro$^\dagger$ (1M) & 21.1 & 46.7 & 42.1 & 44.2 \\
claude-3-sonnet$^\dagger$ (200K) & 21.3 & 45.6 & 40.8 & 43.1 \\
gpt-4-turbo$^\dagger$ (128K) & 21.6 & 51.9 & 46.5 & 48.9 \\
\midrule
\methodname{}（本文） & \ourbest & \ourbest & \ourbest & \ourbest \\
\bottomrule
\multicolumn{5}{@{}l}{\footnotesize RAG 不适用于事件摘要任务（摘要需全局理解而非局部检索）。} \\
\multicolumn{5}{@{}l}{\footnotesize FactScore F1 为本表主指标，ROUGE-L 为辅助参考。} \\
\end{tabular}
\end{table}

\textbf{表~\ref{tab:exp1_layer2_qa} 和~\ref{tab:exp1_layer2_event} 的论证逻辑。}
TPPM 与 LoCoMo 基线在相同数据集（10 段对话）和相同评估协议（QA F1 + FactScore）下比较，确保了方法间可比性。关键比较点：

\begin{enumerate}
  \item \textbf{TPPM vs.\ RAG-Observation}：两者同为"结构化断言检索"范式——RAG-Obs.\ 将对话转为观察（observations），TPPM 将对话转为心理画像（PPMU）。若 TPPM 的 Overall F1 显著高于 RAG-Obs.\ (43.3)，则说明心理画像的结构化设计优于通用观察抽取。

  \item \textbf{TPPM vs.\ gpt-4-turbo}：gpt-4-turbo 是 LoCoMo 最强基线（QA Overall 51.6, FactScore F1 48.9），但其 128K 上下文窗口仍显著落后人类（87.9）。若 TPPM 基于 7B 模型逼近或超过 gpt-4-turbo 的记忆表现，则说明结构化记忆在效率上大幅优于暴力长上下文。

  \item \textbf{TPPM 在 Adversarial 上的鲁棒性}：LoCoMo 原文发现所有 Long-context LLMs 在 Adversarial 问题上灾难性退化（gpt-4-turbo 仅 15.7\%），主要表现为错误归属说话人。TPPM 的场景条件分支机制能否抑制这一退化，是评估其鲁棒性的关键。
\end{enumerate}

\textbf{分析。} 表~\ref{tab:exp1_layer2_qa} 和~\ref{tab:exp1_layer2_event} 的关键价值在于：如果 TPPM 在 Layer 1（PsyDial 回复质量）上体现的咨询技能优势是真实来自更好的心理画像记忆，那么在 Layer 2（LoCoMo 记忆探针）上也应观察到一致的优势。如果 Layer 1 的 OA 提升显著但 Layer 2 提升微弱，说明 TPPM 的增益主要来自其他因素（如更好的 prompt 设计）而非记忆机制本身——这本身就是重要的诊断信息。

\subsubsection{两层结果的综合解读}

Layer 1 和 Layer 2 构成互补验证：

\begin{itemize}
  \item \textbf{Layer 1 回答：}TPPM 生成的回复在心理健康领域中是否具有更好的咨询技能质量（8 维咨询技能 + OA，GPT-4o）？
  \item \textbf{Layer 2 (LoCoMo) 回答：}结构化记忆在通用长程场景中是否有效？LoCoMo 已有的 11 个基线~\cite{maharana2024locomo} 提供了清晰的参照——即使 gpt-4-turbo (128K) 的 QA Overall 也仅 51.6\%，远低于人类的 87.9\%，说明长期记忆仍有巨大提升空间。
\end{itemize}

Layer 1 与 Layer 2 的组合逻辑构成了记忆来源的诊断框架：
\begin{itemize}
  \item ``Layer 1 OA 提升 + Layer 2 LoCoMo 提升'' $\to$ TPPM 通过结构化的心理画像记忆提供了超越原始历史文本的咨询理解能力，记忆机制是有效的增益来源；
  \item ``Layer 1 OA 提升 + Layer 2 LoCoMo 不提升'' $\to$ 需要审视咨询技能增益来源（可能是 prompt 设计或基座模型的领域知识，非记忆机制）；
  \item ``Layer 1 不提升 + Layer 2 LoCoMo 提升'' $\to$ 记忆有效但未有效转化为回复质量（检索策略需优化）。
\end{itemize}

动态记忆机制的设计最优性则由\textbf{实验二}独立验证：Layer 1（PersonaMem~\cite{jiang2025personamem}）验证通用画像槽位抽取能力；Layer 2 通过参数敏感性分析与替代设计对比，验证各动态机制的设计选择是否最优。各机制的不可替代性由\textbf{实验三}消融实验独立验证。

% ============================================================
\subsection{实验二：心理画像记忆建模与动态机制验证}
\label{sec:exp2}

\textbf{实验目的。} 实验一的双层评估验证了 TPPM 在回复生成中的效果及其与通用记忆能力的关联。实验二进一步验证：(1) 系统是否准确抽取了用户画像槽位（画像槽位抽取的通用能力——Layer 1·PersonaMem）；(2) TPPM 的三个动态记忆机制——状态$\to$特质固化、类型条件衰减、场景条件分支——的设计选择是否最优（Layer 2·参数敏感性分析 + 替代设计对比）。

\subsubsection{Layer 1：画像槽位抽取通用能力验证（PersonaMem）}

\textbf{评估协议。}
采用与~\cite{jiang2025personamem} 完全一致的评估协议：在 PersonaMem 的\textbf{therapy consultation 话题子集}（占 PersonaMem 全量数据 $\sim$9.1\%）上，评测 5 类核心查询的 4 选 1 准确率。TPPM 在 Qwen3.5-9B 原生上下文窗口内参与评测——PersonaMem 提供 10/20/60 session 三种长度（$\sim$32K/128K/1M tokens），TPPM 选取上下文窗口内的 session 数量，不依赖长上下文扩展；通过结构化心理画像记忆在有限上下文内压缩跨 session 的关键心理信息，验证记忆效率优势而非上下文容量优势。5 类查询与 TPPM 的画像机制直接对应：

\begin{itemize}
  \item \textbf{查询 1（回忆用户分享的事实）} $\leftrightarrow$ 静态槽位抽取（Profile Slot F1）：模型能否准确回忆用户曾分享的固定信息？
  \item \textbf{查询 3（识别最新偏好）} $\leftrightarrow$ 衰减机制 + 时序一致性（Temporal Consistency）：模型能否基于最近的交互更新对用户的认知？
  \item \textbf{查询 4（追踪完整偏好演化）} $\leftrightarrow$ 状态$\to$特质固化（Trait Consolidation F1）：模型能否追踪用户偏好在多 session 间的变化轨迹？
  \item \textbf{查询 5（回顾偏好变化的原因）} $\leftrightarrow$ 证据集合 $E_i$ + 因果归因：模型能否记住偏好变化背后的具体事件和原因？
  \item \textbf{查询 7（泛化到新场景）} $\leftrightarrow$ 跨场景泛化：模型能否将已学到的画像信息迁移到未知场景？
\end{itemize}

\textbf{设计优势。}
本评估协议具有以下优势：(1) \textbf{与已有基准的可比性}：PersonaMem 已评测 15 个 LLM 基线（GPT-4.5/4.1/4o, Gemini-1.5/2.0-Flash, Claude-3.7/3.5, o4-mini/o3-mini/o1, DeepSeek-R1, Llama-4-Maverick 等）~\cite{jiang2025personamem}；(2) \textbf{领域对齐}：选取 therapy consultation 子集使评测与 TPPM 的心理健康支持场景天然对齐——therapy 话题中用户表达的偏好演化（如对治疗态度、应对策略的变化）与 TPPM 追踪的心理状态变化具有结构相似性；(3) \textbf{效率验证}：TPPM 在 Qwen3.5-9B 原生上下文窗口内参与评测，通过结构化心理画像记忆在有限上下文中压缩关键信息，与直接将全量历史输入长上下文窗口的基线形成效率对比；(4) \textbf{人类参照}：PersonaMem 的 4 选 1 格式中随机基线为 25\%，但最优模型仅 $\sim$52\%，说明长程画像追踪对当前 LLM 仍是未解决问题。

\textbf{局限性说明。} (1) PersonaMem 追踪的是通用偏好/事实（饮食、爱好、职业信息），而非心理状态——即使在 therapy 子集中，画像追踪也与心理状态追踪存在本质差异。因此 PersonaMem 的结果仅作为\textbf{结构化画像抽取的通用能力验证}，不能直接等同于心理画像记忆能力，动态机制的设计最优性由 Layer 2（参数敏感性分析 + 替代设计对比）独立验证。(2) TPPM 仅在 therapy 子集上评测，而引用的 15 个基线数据为 PersonaMem 原文全量结果（15 话题混合）。子集差异导致 TPPM 与基线之间在话题分布上不完全可比——therapy 子集的难度可能系统性地偏离全量平均（如 therapy 对话涉及更多隐式偏好表达），此限制在讨论中进一步说明。(3) TPPM 的跨话题泛化能力未在本实验中验证，由实验一 Layer 2（LoCoMo 10 段对话的通用记忆）提供交叉证据。

\begin{table}[htbp]
\centering
\scriptsize
\setlength{\tabcolsep}{2.3pt}
\caption{实验二·Layer 1：PersonaMem 画像槽位抽取评估（4 选 1 准确率）。\\
评估协议与~\cite{jiang2025personamem} 完全对齐（128K 上下文档位，discriminative setting）。\\
\textbf{x} 为本文方法最优；$<$x 表示劣于本文方法。}
\label{tab:exp2_layer1_personamem}
\begin{tabular}{@{}lcccccc@{}}
\toprule
\multirow{2}{*}{方法} & \multicolumn{6}{c}{PersonaMem 查询类型 (Accuracy) $\uparrow$} \\
\cmidrule(lr){2-7}
& \tabincell{c}{Q1: Recall \\ Facts} & \tabincell{c}{Q3: Latest \\ Pref.} & \tabincell{c}{Q4: Pref. \\ Evolution} & \tabincell{c}{Q5: Reasons \\ for Change} & \tabincell{c}{Q7: Gen- \\ eralize} & Overall \\
\midrule
\multicolumn{7}{c}{\textit{PersonaMem 已评测基线~\cite{jiang2025personamem}}} \\
\midrule
\multicolumn{7}{c}{\textit{Claude 系列}} \\
Claude-3.7$^\dagger$     & 0.25 & 0.09 & 0.45 & 0.57 & 0.29 & 0.26 \\
Claude-3.5$^\dagger$     & 0.29 & 0.27 & 0.55 & 0.64 & 0.20 & 0.30 \\
\midrule
\multicolumn{7}{c}{\textit{GPT 系列}} \\
GPT-4.5$^\dagger$        & 0.61 & 0.55 & 0.68 & 0.76 & 0.46 & 0.52 \\
GPT-4.1$^\dagger$        & 0.65 & 0.50 & 0.67 & 0.84 & 0.53 & 0.52 \\
GPT-4o$^\dagger$         & 0.41 & 0.46 & 0.68 & 0.77 & 0.32 & 0.45 \\
\midrule
\multicolumn{7}{c}{\textit{Gemini 系列}} \\
Gemini-1.5$^\dagger$     & 0.54 & 0.59 & 0.65 & 0.77 & 0.54 & 0.52 \\
Gemini-2.0-Flash$^\dagger$ & 0.50 & 0.52 & 0.70 & 0.79 & 0.46 & 0.49 \\
\midrule
\multicolumn{7}{c}{\textit{推理模型}} \\
o4-mini$^\dagger$        & 0.42 & 0.55 & 0.73 & 0.75 & 0.38 & 0.48 \\
DeepSeek-R1$^\dagger$    & 0.43 & 0.42 & 0.68 & 0.83 & 0.38 & 0.45 \\
\midrule
\multicolumn{7}{c}{\textit{开源模型}} \\
Llama-4-Maverick$^\dagger$ & 0.37 & 0.43 & 0.66 & 0.76 & 0.32 & 0.43 \\
\midrule
\methodname{}（本文） & \ourbest & \ourbest & \ourbest & \ourbest & \ourbest & \ourbest \\
\bottomrule
\multicolumn{7}{@{}l}{\footnotesize 随机基线（4 选 1）= 0.25。Q3=识别最新偏好，Q4=追踪偏好演化，Q5=回顾偏好变化原因，Q7=泛化到新场景。} \\
\multicolumn{7}{@{}l}{\footnotesize TPPM 仅评测 therapy consultation 子集（Qwen3.5-9B 原生上下文窗口内），PersonaMem 基线为原文全量结果~\cite{jiang2025personamem}。} \\
\multicolumn{7}{@{}l}{\footnotesize TPPM 仅评测 therapy consultation 子集（Qwen3.5-9B 原生上下文窗口内），基线数据为 PersonaMem 原文全量结果。} \\
\multicolumn{7}{@{}l}{\footnotesize 子集差异对可比性的影响在讨论中说明；TPPM 的跨领域泛化由实验一 Layer 2（LoCoMo）交叉验证。} \\
\end{tabular}
\end{table}

\textbf{表~\ref{tab:exp2_layer1_personamem} 的论证逻辑。}
TPPM 与 PersonaMem 基线在相同评估协议下比较，确保了指标定义的可比性。TPPM 仅在 therapy consultation 子集上评测，基线数据为 PersonaMem 原文全量结果——话题分布的差异在讨论中作为限制条件说明。关键比较点：(1) 若 Q4（偏好演化）的提升显著大于 Q1（静态回忆），则说明 TPPM 的动态追踪机制（衰减 + 固化）是增益的主要来源；(2) 若 Q5（变化原因回顾）提升显著，则说明 TPPM 的证据集合 $E_i$ 机制有效保留了偏好变化的因果归因信息；(3) 若 Q7（泛化）提升显著，则说明 TPPM 的结构化槽位设计确实支持跨场景迁移，而非仅记忆表面文本。

\subsubsection{Layer 2：动态机制参数敏感性分析}

\textbf{设计动机。}
Layer 1 验证了画像槽位的静态抽取质量，但 TPPM 的核心创新在于动态记忆管理——时间衰减、状态$\to$特质固化、场景条件分支。Layer 2 通过参数敏感性分析与替代设计对比，回答``每个动态机制的设计选择是否最优''。与实验三的消融实验（回答``每个机制是否不可替代''）互补，二者无冗余。

\textbf{2a. 固化机制参数敏感性分析。}
TPPM 的公式 (11) 要求瞬时心理画像记忆在至少 $K$ 个会话中被观测到且固化分数 $\Pi_i(t) > \delta_{\text{promote}}$ 后方可固化为长期画像。我们在 $K \in \{1, 2, 3, 5, 7, 10\}$ 范围内扫描固化阈值，在 PersonaMem 和 LoCoMo 上分别报告 PersonaMem Q4（偏好演化追踪）、Trait Consolidation Precision / Recall 和 LoCoMo QA Multi-hop F1 随 $K$ 的变化。预期存在最优 $K$ 值且呈现清晰的 Precision--Recall 权衡：$K$ 过小导致一次性情绪波动被误固化（Precision 下降），$K$ 过大导致稳定特质因观测窗口不足而无法固化（Recall 下降）。

为验证 TPPM 的 $\Pi$-score 固化公式是否优于简单替代方案，我们在相同 $K$ 值下比较三种固化策略：(1) \textbf{TPPM 固化}（公式 11，$\Pi_i(t) > \delta_{\text{promote}} \wedge N^{\text{session}}_i \geq K$）；(2) \textbf{简单计数固化}（$N^{\text{session}}_i \geq K$，无 $\Pi$-score 门控，纯计数）；(3) \textbf{固定窗口固化}（最近 $K$ 个 session 中至少出现 $K_{\min}$ 次）。若 TPPM 固化优于简单计数固化，则证明 $\Pi$-score 门控不仅``有用''而且``最优''。

\textbf{2b. 衰减机制参数敏感性分析。}
TPPM 的公式 (13) 设定了类型条件衰减率 $\lambda_{\text{affect}} > \lambda_{\text{stressor}} > \lambda_{\text{coping}} \approx \lambda_{\text{support}} > \lambda_{\text{trait}}$。我们在 $\lambda_{\text{scale}} \in \{0.1, 0.3, 0.5, 1.0, 2.0\}$（乘以公式 13 的基础 $\lambda$ 值）范围内扫描衰减率尺度，在 LoCoMo 和 PersonaMem 上分别报告 Temporal F1 和 Q3（识别最新偏好）随 $\lambda_{\text{scale}}$ 的变化，以及各维度 Per-Slot F1 对衰减率的敏感性差异。

为验证差分衰减是否优于替代方案，我们比较四种衰减策略：(1) 无衰减（$\lambda_d \equiv 0$）；(2) 统一衰减（单一 $\lambda$，所有维度相同）；(3) 随机差分衰减（维度间有差异但非 TPPM 的心理学排序）；(4) TPPM 差分衰减（$\lambda_{\text{affect}} > \lambda_{\text{stressor}} > \lambda_{\text{coping}} \approx \lambda_{\text{support}} > \lambda_{\text{trait}}$）。关键比较：若 TPPM 差分衰减 $>$ 随机差分衰减 $>$ 统一衰减 $>$ 无衰减，则证明``不仅需要衰减，还需要按心理学排序的差分衰减''——这是最强的论证。若 TPPM 差分衰减 $\approx$ 随机差分衰减 $>$ 统一衰减，则仅证明``需要差分但排序不重要''，需在讨论中承认。

\textbf{2c. 场景分支机制参数敏感性分析。}
TPPM 的公式 (3) 通过上下文相似度阈值 $\delta_{\text{ctx}}$ 控制何时为心理画像建立场景条件分支。我们在 $\delta_{\text{ctx}} \in \{0.3, 0.5, 0.7, 0.9\}$ 范围内扫描分支阈值，在 PersonaMem 和 LoCoMo 上分别报告 Q7（跨场景泛化）和 Adversarial F1 随 $\delta_{\text{ctx}}$ 的变化，以及分支数量随 $\delta_{\text{ctx}}$ 的变化。

为验证连续分支是否优于粗粒度替代方案，我们比较三种分支策略：(1) 无分支（全局画像更新，$\delta_{\text{ctx}} \to \infty$）；(2) 二值分支（仅区分``危机''vs.``非危机''两种场景）；(3) TPPM 连续分支（公式 3，基于上下文相似度的软分支）。若连续分支 $>$ 二值分支 $>$ 无分支，则场景分支有效且细粒度分支优于粗粒度。

\textbf{Layer 2 论证逻辑总结。}
Layer 2 的三个参数敏感性分析回答``每个动态机制的设计选择是否最优''：
\begin{itemize}
  \item 固化敏感性：$K$ 最优值存在 $+$ $\Pi$-score 优于简单计数 $\to$ 固化设计不仅``有用''而且``最优''
  \item 衰减敏感性：$\lambda$ 最优值存在 $+$ 心理学排序差分优于随机差分 $\to$ 衰减排序有理论依据
  \item 分支敏感性：$\delta_{\text{ctx}}$ 最优值存在 $+$ 连续分支优于二值分支 $\to$ 分支粒度有实际意义
\end{itemize}

与实验三消融（``是否不可替代''）互补，形成完整论证：实验二 Layer 1 证明画像更准，Layer 2 证明设计最优，实验三证明不可替代。

\subsubsection{画像建模指标的人工验证}
\label{sec:exp2_human_val}

\textbf{验证动机。} 自动化指标可能无法完全反映画像质量的真实感知。为此，我们设计一项小规模的人工验证研究，检验自动化画像指标与专家判断的一致性。

\textbf{验证设计。} 从 PersonaMem therapy 子集和 PsyDial 测试集中各随机抽取 15 段对话（共 30 段），由 2 名具有心理咨询资质的专家（与研究团队无关的外部专家，具有 5 年以上临床经验）独立进行以下两项判断：
\begin{enumerate}
  \item \textbf{画像充分性评分（Profile Adequacy Rating）}：对系统输出的心理画像进行 1--5 Likert 评分——``该画像是否准确、完整地反映了用户在指定会话区间内的心理状态？''
  \item \textbf{固化恰当性判断（Consolidation Appropriateness）}：对系统标记为``已固化''的每个长期特质，判定为``应固化''或``不应固化''（二分判断）。
\end{enumerate}
专家之间的一致性使用 Krippendorff's $\alpha$ 衡量（目标 $\alpha \geq 0.75$）。效度证据通过计算自动化画像指标（Profile Slot F1、Trait Consolidation Precision/Recall）与专家评分的相关性获得。

\subsubsection{跨实验的逐样本关联分析}

\textbf{分析动机。} 实验一和实验二分别从``生成质量''和``记忆准确性''两个侧面评估 TPPM，但二者之间的关联仍是推断性的。如果 TPPM 的共情质量优势确实来源于更好的心理画像记忆，则应能在逐样本层面观察到：\textbf{记忆准确性更高的样本，其生成质量也应更高}。

\textbf{实验设计。} 对 PsyDial 测试集中每个评估样本（N $\approx$ 100），同时获得其实验一的 OA 总分和实验二中的 Profile Slot F1。计算二者的 Pearson $r$ 和 Spearman $\rho$。预期 \methodname{} 条件下 $r > 0$（正相关），表明系统确实利用了准确的画像信息来生成更好的回复。如果 \methodname{} 条件下关联显著，则提供了比分组趋势更强的机制证据。

\subsubsection{按心理维度的画像槽位细粒度分析}

TPPM 将心理画像属性槽位 $a_i$ 设计为七个心理相关子空间~\cite{engel1977biopsychosocial,trzepacz1993mse}：情绪状态（affective state）、压力来源（stressor）、认知表征（cognitive pattern）、应对方式（coping style）、社会支持需求（support need）、行为与生活节律（behavioral routine）以及风险信号（risk signal）。对每个心理维度 $d \in \mathcal{A}$，我们额外计算该维度上的 Per-Slot F1$^{(d)}$，以诊断不同维度的抽取难度差异。预期认知表征和风险信号等内隐维度的 F1 系统性地低于情绪状态和压力来源等显式维度，而 TPPM 的结构化槽位设计在隐式维度上的相对增益更大。

\subsubsection{会话数量的敏感性分析}

\textbf{分析动机。} 此处进一步检验：随着会话数量增长，记忆质量如何变化？

我们从 PsyDial 测试集中按会话数量（$N=1,3,5,10,15,20$）分层抽样 6 个子测试集（各 100 段），测量 Profile Slot F1、Trait Consolidation F1 和 Unsupported Profile Rate 随 $N$ 的变化。预期 Long-Context 随 $N$ 增长趋于饱和甚至退化，而 TPPM 随 $N$ 增长持续受益于证据积累——两条曲线的趋势差是长期记忆能力的最有力证据。

\textbf{抽样混杂控制。} 由于 PsyDial 中 $N \geq 15$ 的长对话可能在对话主题、来访者类型和问题复杂度上与短对话存在系统性差异（长对话更可能涉及慢性心理困扰或深度咨询关系），高 $N$ 子集和低 $N$ 子集之间的差异可能掺杂对话类型的混杂效应。为部分分离会话数量效应与对话类型效应，我们在报告完整分层结果的同时，额外进行一项\textbf{窗口内比较（within-window comparison）}：对同一段长对话（$N \geq 10$），取其前 $k$ 个会话（$k \in \{3, 5, 7, 10\}$）作为子序列，测量随 $k$ 增长的性能变化。由于所有子序列来自同一对话，对话主题和来访者类型被自然控制，剩余的变化更可能归因于会话数量本身。

\textbf{补充：记忆时长—遗忘率曲线。} 在 PersonaMem 上，测量系统对同一心理信号的保持率随距首次观测的时间间隔的变化。预期 \methodname{} 的遗忘曲线显著浅于无衰减条件。

% ============================================================
\subsection{实验三：消融实验}
\label{sec:exp3}

\textbf{实验目的。} 分析 TPPM 五个核心模块各自对最终性能的贡献。我们依次移除状态$\to$特质固化、场景条件分支、时间衰减、结构化心理槽位以及长期画像检索，在实验一和实验二的全部关键指标上报告性能退化。此外，增设扰码记忆和噪声记忆两项对照条件，以隔离``格式效应''与``内容效应''。

\subsubsection{消融设计与指标覆盖}

\textbf{五个消融变体。}
\begin{enumerate}
  \item \textbf{移除状态$\to$特质固化（w/o Consolidation）}：关闭公式 (11) 的固化门控——所有瞬时画像记忆 $m_i^{\text{ep}}$ 在会话结束后直接丢弃，不再积累为长期特质 $\mathcal{T}$。预期主要影响：Trait Consolidation F1 退化最大，长期画像缺失导致 PsyDial OA 和 LoCoMo QA F1 下降。
  \item \textbf{移除场景条件分支（w/o Scene-Conditional Branching）}：将公式 (3) 的场景条件分支替换为无条件的全局画像更新——系统不再根据对话场景（初次评估 vs.\ 危机干预 vs.\ 日常随访）切换画像读取/写入策略。预期主要影响：PersonaMem Q7（跨场景泛化）和 LoCoMo Adversarial F1 退化（无法区分情境差异与真实心理转折），矛盾信号压力测试中的 False Conflict Rate 上升。
  \item \textbf{移除时间衰减（w/o Temporal Decay）}：关闭公式 (13) 的类型条件衰减——所有心理画像属性值 $v(a_i)$ 在时间推移中保持不变（即设置 $\lambda_d \equiv 0, \forall d$）。预期主要影响：陈旧信息干扰导致 LoCoMo QA Temporal F1 和 PersonaMem Q3（最新偏好识别）退化，Profile Slot F1 中的时序一致性下降。
  \item \textbf{移除结构化心理槽位（w/o Structured Slots）}：将 7 维 PPMU 结构化槽位替换为自由文本的心理描述段落（等长但无固定槽位结构）。预期主要影响：Unsupported Profile Rate 恶化最大（生成式段落的幻觉无槽位约束），Profile Slot F1 下降，PersonaMem Q7（泛化）退化。
  \item \textbf{移除长期画像检索（w/o Long-term Retrieval）}：在生成回复时仅使用当前会话的瞬时记忆 $m_i^{\text{ep}}$，不从记忆 bank 中检索已固化的长期特质 $\mathcal{T}$ 和历史瞬时记忆。预期主要影响：PsyDial OA 和 LoCoMo QA Overall F1 退化最大（丢失跨会话信息），PersonaMem Q4（偏好演化追踪）退化。
\end{enumerate}

\textbf{指标选择。} 消融实验覆盖实验一和实验二的 4 个关键指标，确保每个模块的消融效应可在对应维度上被诊断：

\begin{itemize}
  \item \textbf{实验一·Layer 1（PsyDial）}：GPT-4o OA $\uparrow$ — 下游回复质量
  \item \textbf{实验一·Layer 2（LoCoMo）}：QA Overall F1 $\uparrow$ — 通用记忆保持
  \item \textbf{实验二·Layer 1（PersonaMem）}：Overall Accuracy $\uparrow$ — 画像槽位抽取的端到端行为准确率
  \item \textbf{画像质量}：Unsupported Profile Rate $\downarrow$ — 幻觉控制
\end{itemize}

此外，对 PersonaMem 拆出 Q3（最新偏好识别）、Q4（偏好演化追踪）、Q7（跨场景泛化）子指标，对 LoCoMo 拆出 Temporal F1 和 Adversarial F1 子指标，以提供更细粒度的消融诊断。

\subsubsection{消融结果}

\begin{table}[htbp]
\centering
\footnotesize
\setlength{\tabcolsep}{2.5pt}
\caption{实验三·消融实验：逐一移除核心模块后的性能退化。\\
\textbf{加粗}的 $\Delta$ 表示预期退化幅度最大的指标—模块对。$\Delta_i^{\text{(+)}}>0$ 表示指标下降（恶化），$\Delta_i^{\text{(-)}}>0$ 表示指标上升（恶化，适用于 Unsupported Rate）。}
\label{tab:ablation}
\begin{tabular}{@{}lccccc@{}}
\toprule
\multirow{2}{*}{模型变体} & \multicolumn{2}{c}{实验一} & 实验二·L1 & \multirow{2}{*}{\tabincell{c}{Unsup. \\ Rate $\downarrow$}} & \multirow{2}{*}{\tabincell{c}{子指标诊断 \\ （最大退化）}} \\
\cmidrule(lr){2-3} \cmidrule(lr){4-4}
& \tabincell{c}{PsyDial \\ OA $\uparrow$} & \tabincell{c}{LoCoMo \\ QA F1 $\uparrow$}
& \tabincell{c}{PersonaMem \\ Acc. $\uparrow$} & & \\
\midrule
\methodname{}（完整） & \ourbest & \ourbest & \ourbest & \ourbest & --- \\
\midrule
\multicolumn{6}{c}{\textit{核心模块消融}} \\
\midrule
w/o 状态$\to$特质固化
  & $x-\Delta_1^{\text{(+)}}$ & $x-\Delta_2^{\text{(+)}}$ & $x-\Delta_3^{\text{(+)}}$ & $x+\Delta_4^{\text{(-)}}$ & \textbf{PMem Q4 $\downarrow$, LoCoMo Multi-hop $\downarrow$} \\[4pt]
w/o 场景条件分支
  & $x-\Delta_5^{\text{(+)}}$ & $x-\Delta_6^{\text{(+)}}$ & $x-\Delta_7^{\text{(+)}}$ & $x+\Delta_8^{\text{(-)}}$ & \textbf{PMem Q7 $\downarrow$, LoCoMo Adversarial $\downarrow$} \\[4pt]
w/o 时间衰减
  & $x-\Delta_9^{\text{(+)}}$ & \textbf{$x-\Delta_{10}^{\text{(+)}}$} & $x-\Delta_{11}^{\text{(+)}}$ & $x+\Delta_{12}^{\text{(-)}}$ & \textbf{LoCoMo Temporal $\downarrow$, PMem Q3 $\downarrow$} \\[4pt]
w/o 结构化心理槽位
  & $x-\Delta_{13}^{\text{(+)}}$ & $x-\Delta_{14}^{\text{(+)}}$ & \textbf{$x-\Delta_{15}^{\text{(+)}}$} & \textbf{$x+\Delta_{16}^{\text{(-)}}$} & PMem Q7 $\downarrow$ \\[4pt]
w/o 长期画像检索
  & \textbf{$x-\Delta_{17}^{\text{(+)}}$} & \textbf{$x-\Delta_{18}^{\text{(+)}}$} & \textbf{$x-\Delta_{19}^{\text{(+)}}$} & $x+\Delta_{20}^{\text{(-)}}$ & PMem Q4 $\downarrow$ \\
\midrule
\multicolumn{6}{c}{\textit{对照条件}} \\
\midrule
+ 扰码记忆（格式对照）
  & $x-\Delta_{21}^{\text{(+)}}$ & $x-\Delta_{22}^{\text{(+)}}$ & \multicolumn{3}{c}{N/A（记忆非真实画像，无对应 GT）} \\[4pt]
+ 噪声记忆（20\% 错误）
  & $x-\Delta_{23}^{\text{(+)}}$ & $x-\Delta_{24}^{\text{(+)}}$ & $x-\Delta_{25}^{\text{(+)}}$ & $x+\Delta_{26}^{\text{(-)}}$ & --- \\
\bottomrule
\end{tabular}

\vspace{4pt}
{\footnotesize 注：(1) $\Delta_i>0$；(2) \textbf{加粗}表示该消融条件下预期退化幅度最大的指标——验证每个模块的独立功能；(3) 子指标诊断列报告各消融条件下退化幅度最大的细粒度指标（PersonaMem Q3/Q4/Q7、LoCoMo Temporal/Adversarial/Multi-hop），提供主指标之外的诊断信息；(4) 扰码记忆对照的 PersonaMem 指标 N/A 因为扰码画像无对应正确 ground truth 用于 4 选 1 判断。}
\end{table}

\subsubsection{模块—指标退化对应分析}

表~\ref{tab:ablation} 预期揭示清晰的模块—指标退化对应模式，为每个 TPPM 组件提供独立的功能验证：

\begin{enumerate}
  \item \textbf{状态$\to$特质固化 $\to$ PersonaMem Q4 / LoCoMo Multi-hop F1}：固化机制是将跨会话的重复心理信号积累为长期特质的关键。移除后，PersonaMem Q4（偏好演化追踪）和 LoCoMo Multi-hop F1 退化最大——系统失去了区分``一次性情绪波动''与``稳定心理特质''的能力，无法有效积累跨会话信息。Trait Consolidation Precision/Recall 预期大幅下降。

  \item \textbf{场景条件分支 $\to$ PersonaMem Q7 / LoCoMo Adversarial F1}：场景条件分支使系统能够根据对话情境（初次评估/危机干预/日常随访）切换画像更新策略。移除后，PersonaMem Q7（跨场景泛化）和 LoCoMo Adversarial F1 退化最大——系统将情境差异误判为真实心理转折（False Conflict Rate 上升）。这也将在对抗性压力测试的矛盾信号测试中得到进一步验证。

  \item \textbf{时间衰减 $\to$ LoCoMo QA Temporal F1 / PersonaMem Q3}：类型条件衰减率（$\lambda_{\text{affect}} > \lambda_{\text{stressor}} > \lambda_{\text{trait}}$）使系统对不同心理维度施加差异化的遗忘速率。移除后（统一不衰减），陈旧信息干扰当前判断——LoCoMo QA 的 Temporal F1（时间推理）和 PersonaMem Q3（识别最新偏好）退化最为明显。时间衰减在长程场景中的价值更大，实验二 Layer 2 的衰减策略对比进一步验证了差分衰减的优越性。

  \item \textbf{结构化心理槽位 $\to$ PersonaMem Accuracy / Unsupported Rate}：7 维 PPMU 结构化槽位约束了画像输出的格式和内容边界。移除后，自由文本段落的幻觉无法被槽位边界抑制——Unsupported Profile Rate 恶化最大，PersonaMem Overall Accuracy（受幻觉影响直接降低 4 选 1 准确率）退化显著。此外，PersonaMem Q7（跨场景泛化）退化幅度预期大于 Q1（事实回忆），因为结构化槽位的类型约束对泛化推理的支持大于对静态事实回忆的支持。

  \item \textbf{长期画像检索 $\to$ PsyDial OA / LoCoMo QA F1 / PersonaMem Q4}：长期画像检索是连接记忆 bank 与生成器的桥梁——将固化特质 $\mathcal{T}$ 和历史瞬时记忆注入回复生成的上下文。移除后，系统退化为仅依赖当前会话瞬时记忆的``短程''系统——PsyDial OA 和 LoCoMo QA Overall F1 退化最大（失去跨会话信息），PersonaMem Q4（偏好演化追踪，需要跨 session 检索历史画像）退化显著。
\end{enumerate}

\textbf{跨实验的消融退化模式诊断。} 表~\ref{tab:ablation} 的跨实验指标覆盖使我们能够进一步区分``记忆贡献''与``非记忆贡献''：

\begin{itemize}
  \item 如果移除长期画像检索后 PsyDial OA 显著下降但 LoCoMo QA F1 变化不大，说明画像检索的增益并非通过改善记忆保持实现（可能是 prompt 格式效应），需要重新审视扰码记忆对照的结果；
  \item 如果移除时间衰减后 PersonaMem Q3 退化但 PsyDial OA 不退化，说明时间衰减对短程对话（PsyDial 平均 37.8 轮/对话，时间跨度有限）的贡献有限，其价值主要体现在 LoCoMo 等长程场景中；
  \item 如果移除结构化心理槽位后 Unsupported Rate 上升但 PersonaMem Accuracy 不降，说明幻觉增加并未影响端到端行为（系统可能通过其他线索补偿了错误画像信息），这是一种需要警惕的``伪鲁棒性''。
\end{itemize}

\subsubsection{对照条件的解释逻辑}

\textbf{扰码记忆对照（Scrambled Memory Control）。} 该对照保留完全相同的结构化心理画像格式（7 维 PPMU，含槽位名、属性值、证据来源），但将画像内容替换为随机生成的、语法正确但内容错误的伪画像信息（如随机交换不同来访者的画像槽位值）。其论证逻辑如下：

\begin{itemize}
  \item 若完整 TPPM $\gg$ 扰码记忆 $\approx$ w/o 结构化槽位：正确的记忆内容是驱动性能提升的核心因素，TPPM 的机制验证逻辑成立；
  \item 若完整 TPPM $\approx$ 扰码记忆 $>$ w/o 结构化槽位：prompt 中``存在心理画像格式段落''这一事实（而非其正确内容）是增益来源——核心创新论点被实质性削弱，需重新审视结构化槽位设计的有效性；
  \item 若完整 TPPM $>$ 扰码记忆 $>$ w/o 结构化槽位：格式和内容各自贡献了一部分增益（混合效应），需进一步分析各自贡献比例。
\end{itemize}

\textbf{噪声记忆对照（Noisy Memory Control）。} 在 TPPM 记忆 bank 中随机插入 20\% 的错误画像信息（如故意将``焦虑水平：显著下降''改写为``焦虑水平：显著上升''），检验系统对记忆质量的敏感性：

\begin{itemize}
  \item 若噪声记忆导致 PsyDial OA 和 PersonaMem Accuracy 显著退化：系统确实依赖记忆的具体内容（而非仅利用其``存在感''），鲁棒性验证通过；
  \item 若噪声记忆条件下性能未显著退化：系统可能仅利用了记忆的``存在感''（如心理画像的存在提升了对用户的关注度，但未实际利用其内容）——这是一种需要警惕的虚假鲁棒性，可能意味着模型主要依赖对话表面的情感线索而非画像内容生成回复。
\end{itemize}

\textbf{噪声比例的梯度设计。} 为刻画系统对记忆质量的敏感度曲线，除 20\% 基础条件外，在 PsyDial OA 上额外测试 0\%（完整 TPPM 上界）、10\%、30\%、50\% 四个噪声比例，报告 OA 随噪声比例下降的趋势。预期 OA 随噪声比例接近线性下降——若在 30\% 噪声下 OA 仍不退化，则存在较强的虚假鲁棒性。

\subsubsection{与实验一、实验二的逻辑闭环}

\textbf{消融实验在三实验框架中的角色。} 实验一验证了 TPPM 在回复质量（PsyDial）和通用记忆保持（LoCoMo）上的效果——回答了\textbf{``是否更好''}。实验二验证了这些效果背后的机制——画像槽位抽取能力（PersonaMem）和各动态机制的设计最优性（参数敏感性分析 + 替代设计对比）——回答了\textbf{``为何更好''}。实验三的消融则闭合了最后一个环节——\textbf{``各组件是否不可替代''}：

\begin{itemize}
  \item 若某模块的移除导致实验一中回复质量下降但实验二中记忆指标不变，则说明该模块的贡献路径独立于记忆机制（可能通过 prompt 优化等非记忆渠道）；
  \item 若某模块的移除导致实验二中记忆指标下降但实验一中回复质量不变，则说明记忆精度提升未有效转化为回复质量——检索策略或生成器对记忆信号的利用效率需要改进；
  \item 若某模块的移除同时导致实验一和实验二的指标退化，且退化模式与模块的设计功能一致（如固化模块 $\to$ PersonaMem Q4 + Trait Cons.\ F1），则为该模块提供了\textbf{跨实验的收敛验证}。
\end{itemize}

\textbf{扰码记忆对照在因果链中的关键作用。} 扰码记忆对照是区分``格式效应''与``内容效应''的唯一实验设计——它直接检验实验一中 TPPM 的 OA 优势是否来源于正确的画像内容（而非 prompt 中出现了结构化画像这一格式）。若扰码记忆对照的结果支持``内容效应''解释，则实验二（画像质量验证）和实验三（模块消融）的结果获得了因果解释的坚实基础；反之，则需要对整套论证逻辑进行根本性修正。

\subsection{对抗性压力测试}
\label{sec:stress}

\textbf{设计动机。} 所有上述实验均为验证性（confirmatory）测试——设计用于\textbf{确认} TPPM 的优势。为全面评估方法的鲁棒性，我们进一步设计三类\textbf{对抗性压力测试}——专门用于\textbf{揭示} TPPM 的失败模式。

\begin{enumerate}
  \item \textbf{矛盾信号压力测试（Contradictory Signal Stress Test）}：构造 $N=30$ 段对话（基于功效分析：检测中等效应量 $d=0.5$ 时 $\alpha=0.05$、power $\geq 0.80$），其中来访者在不同会话中表达截然相反的心理信号。\textbf{量化指标：}(a) \textbf{矛盾检测率（Contradiction Detection Rate）}：系统是否标记了矛盾信号的存在（二分判断）；(b) \textbf{情境归因准确率（Context Attribution Accuracy）}：对于检测到的矛盾，系统是否通过场景条件分支正确地将表面矛盾归因于情境差异（而非错误地触发心理转折标记）；(c) \textbf{False Conflict Rate}：系统将情境差异误判为真实心理转折的比例。
  \item \textbf{高度不稳定心理状态测试（High-Instability Test）}：从 PsyDial 中按跨会话情绪变化方差筛选 $N=30$ 段最高波动的对话（由 2 名标注员独立评定方差等级，取交集确认）。\textbf{量化指标：}(a) \textbf{固化抑制率（Consolidation Suppression Rate）}：当 $K=5$ 时，高度不稳定用户中\textbf{不应固化}的瞬时记忆实际被固化的比例——预期应接近 0；(b) \textbf{波动-特质误判率（Fluctuation-Trait Misclassification Rate）}：系统将高度波动的短期心理信号标记为长期特质的比例。
  \item \textbf{心理画像模糊测试（Profile Ambiguity Test）}：构造 $N=30$ 段对话（同功效分析），其中来访者的心理状态描述保持模糊。\textbf{量化指标：}(a) \textbf{Precision@Top-K}：当信号不明确时，系统输出的心理画像中正确抽取的比例；(b) \textbf{Unsupported Profile Rate}：模糊输入下仍被系统写入画像的比例——预期应接近 0。
\end{enumerate}

\textbf{统计性质说明：}三类压力测试的目的不是多方法排序比较（因此不适用主实验的 FDR 校正框架），而是\textbf{定性—定量混合的诊断分析}：量化指标用于测量特定失败模式的严重程度，后续定性案例分析用于理解失败的具体机制。样本量 $N=30$ 足以检测 $d \geq 0.52$ 的效应（paired $t$-test, $\alpha=0.05$, power=0.80），但对于更小的效应量仅具探索性质。压力测试结果在讨论部分与主实验结果合并解读，用于勾勒 TPPM 的失败边界和改进方向。

% ============================================================
\section{讨论}
\label{sec:discussion}

综合三个实验的结果，我们得出以下结论：

\textbf{第一，长程心理支持对话的核心挑战不只是生成能力，更是时间维度上的个体建模能力。}
实验一的双层评估提供了直接证据：Layer 1（PsyDial）验证了 TPPM 在领域回复质量上的咨询技能优势（8 维咨询技能 + OA 显著领先）。Layer 2（LoCoMo）进一步验证了这种优势来源于真实的心理状态记忆保持——当前最强模型（gpt-4-turbo，QA Overall 51.6\%，FactScore F1 48.9）与人类（87.9\%）之间的巨大鸿沟揭示了长程记忆的固有难度，TPPM 的结构化记忆正是瞄准这一鸿沟。实验二的双层设计进一步验证了心理画像记忆的机制：Layer 1（PersonaMem）验证了画像槽位抽取的通用能力，Layer 2 通过参数敏感性分析与替代设计对比，验证了各动态机制设计选择的最优性——固化阈值 $K$ 存在最优值、类型条件差分衰减优于统一衰减和随机差分衰减、连续场景分支优于二值分支。实验三的消融实验则验证了各模块的不可替代性。会话数量敏感性分析佐证：TPPM 随会话增长持续受益于证据积累。

\textbf{第二，心理画像记忆比通用对话记忆更适合心理健康场景。}
与 SoulChat、CPsyCounX、$\pi_4$ 等心理咨询模型相比，TPPM 的优势不仅限于咨询技能评分（Table~\ref{tab:exp1_layer1_auto}），更体现在长期特质固化（实验二）和跨会话记忆保持（Layer 2 LoCoMo）上。LoCoMo 结果说明即使是最强的 RAG-Observation 基线（Overall F1 43.3\%），仍显著落后人类（87.9\%），而 TPPM 通过心理画像的结构化设计有望进一步提升。实验二的 PersonaMem 评估验证了结构化画像抽取的通用性——即使最强模型 GPT-4.5 也仅 $\sim$52\%（4 选 1），而 TPPM 的槽位设计在偏好演化追踪（Q4）和跨场景泛化（Q7）上预期带来显著增益。

\textbf{第三，各模块协同工作，且各自承担明确的职责。}
实验三的消融分析揭示了清晰的模块—指标退化对应模式：固化机制支撑 PersonaMem Q4，场景分支支撑 Q7 和 LoCoMo Adversarial F1，结构化槽位约束幻觉，时间衰减避免陈旧信息干扰。各模块消融均导致预期对应的关键指标退化最大，且退化模式与设计假说一致。

\textbf{第四，方法的有效性不依赖特定基座模型，且在实际部署中可行。}
实验三的消融实验为每个 TPPM 组件提供了独立的功能验证——各模块移除后其预期对应的关键指标退化最大，验证了模块功能假说。扰码记忆和噪声记忆对照为区分``格式效应''与``内容效应''提供了关键证据。

\textbf{第五，心理画像维度的文化效度需要审慎考量。}
TPPM 的七个 PPMU 维度以 Engel（1977）的生物心理社会模型和 Trzepacz \& Baker（1993）的精神状态检查为理论基础，这些框架主要源自西方临床心理学传统。在中文心理咨询语境中（PsyDial），心理健康表达可能存在系统性差异——如躯体化表达（somatization）比情绪化表达更常见，家庭关系而非个体自主性构成压力来源的核心框架~\cite{kirmayer2001cultural}。因此，在报告各维度的 Per-Slot F1 时应按维度分别讨论文化差异可能带来的系统偏差，并在讨论中审慎避免将低 F1 简单归因于``系统能力不足''——低 F1 可能部分反映了该维度在中文语境下的可标注性（annotatability）本身较低。此外，在非中英文语境下的表现需独立验证。

\textbf{第六，TPPM 在极端条件下的失败模式需要被认识和改进。}
对抗性压力测试的目的不是验证方法优势，而是揭示方法的\textbf{失败边界}：(1) 当来访者表达高度矛盾的心理信号时，场景条件分支机制能否通过情境差异解释表面矛盾，而非平滑忽略冲突？(2) 当心理状态极度不稳定时，系统是否会抑制过早固化（大 $K$），还是将波动误判为特质？(3) 当心理信息本身模糊时，系统是否保持低 Unsupported Profile Rate 而非依赖猜测？这些压力测试结果为后续迭代提供了明确的方向，也限制了主实验结果的外推范围。

\textbf{第七，已知局限性与未评估的维度。}
(1) \textbf{心理健康纵向特异性验证的缺失：}本实验使用的三个数据集（PsyDial、LoCoMo、PersonaMem）均非心理健康纵向追踪数据集——PsyDial 评估领域回复质量，LoCoMo 评估通用记忆保持，PersonaMem 评估通用画像追踪。TPPM 的核心假设是心理画像记忆能够捕捉心理状态的动态演化过程，而当前实验无法直接验证这一假设在心理健康纵向数据上的特异性。未来工作应在真实的心理健康纵向数据（如心理咨询跨session追踪、心理健康状态变化时间线）上评估 TPPM 的动态记忆机制。(2) \textbf{因果推断的局限性：}虽然 Layer 1（PsyDial）+ Layer 2（LoCoMo）+ 实验二（画像建模 + 敏感性分析）的多证据链和逐样本关联分析提供了多层次的趋同证据（converging evidence），且扰码记忆对照为隔离格式效应提供了关键诊断，但这仍属于观察性证据而非严格的因果证明。在缺少随机化实验设计（如在相同对话上随机开关记忆机制）的情况下，不能排除未测量的混淆因素。(3) \textbf{患者体验未测量：}所有评估指标均从外部评价者（自动指标、标注员、专家）的角度评估系统输出，但缺乏最关键的利益相关者视角——\textbf{实际心理健康服务使用者的体验}。系统是否能被求助者感知为``被理解''和``被帮助''，需要在未来的用户研究中单独验证。(4) \textbf{治疗师视角缺失：}未评估 TPPM 生成的画像是否对临床实践具有实际价值——治疗师是否信任系统维护的心理画像？是否会将其作为决策参考？治疗师—AI 协作场景（TPPM 辅助而非替代人类治疗师）的评估协议留待未来工作开发。(5)\textbf{安全性评估的局限性：}本文采用 GPT-4o 的 8 维咨询技能评估（含 NJAA 维度）作为质量评判，但未包含专门的安全性评估协议（如风险信号漏检、病理性正常化等子维度）。GPT-4o 的 NJAA 评分可部分反映非评判性态度，但不足以替代面向心理健康场景的专项安全性测试。对于面向脆弱人群部署的系统，应在未来工作中引入专项安全性评估和部署后持续监控。\textbf{跨语言和文化外推}已在第五节中讨论。

% ============================================================
%  第九节：结论
% ============================================================
\section{结论}

本文针对心理健康对话场景中的会话内遗忘与跨会话遗忘问题，提出了时间心理画像记忆（TPPM）框架。TPPM 以心理画像记忆单元（PPMU）为核心表示，在工作记忆、瞬时心理画像记忆和长期心理画像记忆三层结构中组织心理状态信息；通过 LLM 驱动的心理语义抽取与显式记忆演化机制的混合式设计，实现对个体心理状态在时间维度上的细粒度建模。实验结果表明：(1) TPPM 在 PsyDial 和 LoCoMo 上的双层评估中，在咨询技能质量和跨会话记忆保持两个维度均展现出优于现有基线方法的性能；(2) PersonaMem 上的评估验证了结构化心理画像抽取能力，参数敏感性分析进一步验证了各动态机制（固化、衰减、分支）设计选择的最优性；(3) 消融实验为五个核心模块各自的不可替代性提供了独立验证；(4) 对抗性压力测试勾勒了 TPPM 在矛盾信号处理、高度不稳定心理状态和模糊输入条件下的失败边界，为后续改进提供了明确方向。

未来工作将沿以下几个方向展开：第一，在心理健康纵向追踪数据上评估 TPPM 的动态记忆机制，验证其在心理状态变化检测上的特异性；第二，引入随机化实验设计（如在相同对话上随机开关记忆机制），从因果推断层面验证画像记忆机制对生成质量的贡献路径；第三，开展以实际心理健康服务使用者为中心的体验评估，测量用户对系统的感知同理心、被理解感和帮助意愿；第四，探索治疗师—AI 协作场景下的画像共享与临床决策支持协议；第五，强化跨语言、跨文化场景下的心理画像维度效度验证，确保 TPPM 在非中文语境中同样适用且不引入文化偏差；第六，面向实际部署需求，在更长的交互时间尺度（数月乃至数年）上验证画像记忆的长期稳定性和可维护性。

% ============================================================
\bibliographystyle{plainnat}
\begin{thebibliography}{99}

\bibitem{qiu2025psydial}
Qiu, H. \& Lan, Z. PsyDial: A Large-scale Long-term Conversational Dataset for Mental Health Support. \textit{ACL}, 2025.

\bibitem{maharana2024locomo}
Maharana, A. \textit{et al.} Evaluating Very Long-Term Conversational Memory of LLM Agents. \textit{ACL}, 2024.

\bibitem{jiang2025personamem}
Jiang, Y. \textit{et al.} Know Me, Respond to Me: Benchmarking LLMs for Dynamic User Profiling and Personalized Responses at Scale. \textit{COLM}, 2025.

\bibitem{zhang2024cpsycoun}
Zhang, C. \textit{et al.} CPsyCoun: A Report-based Multi-turn Dialogue Reconstruction and Evaluation Framework for Chinese Psychological Counseling. \textit{Findings of ACL}, 2024.

\bibitem{chen2023soulchat}
Chen, Y. \textit{et al.} SoulChat: Improving LLMs' Empathy, Listening, and Comfort Abilities through Fine-tuning with Multi-turn Empathy Conversations. \textit{Findings of EMNLP}, 2023.

\bibitem{qiu2024smile}
Qiu, H. \textit{et al.} SMILE: Single-turn to Multi-turn Inclusive Language Expansion via ChatGPT for Mental Health Support. \textit{Findings of EMNLP}, 2024.

\bibitem{qiu2024psychat}
Qiu, H. \textit{et al.} PsyChat: A Client-centric Dialogue System for Mental Health Support. \textit{CSCWD}, 2024.

\bibitem{zheng2023judging}
Zheng, L. \textit{et al.} Judging LLM-as-a-Judge with MT-Bench and Chatbot Arena. \textit{NeurIPS}, 2023.

\bibitem{engel1977biopsychosocial}
Engel, G. L. The Need for a New Medical Model: A Challenge for Biomedicine. \textit{Science}, 196(4286):129--136, 1977.

\bibitem{trzepacz1993mse}
Trzepacz, P. T. \& Baker, R. W. \textit{The Psychiatric Mental Status Examination}. Oxford University Press, 1993.

\bibitem{kirmayer2001cultural}
Kirmayer, L. J. Cultural Variations in the Clinical Presentation of Depression and Anxiety: Implications for Diagnosis and Treatment. \textit{Journal of Clinical Psychiatry}, 62(Suppl 13):22--30, 2001.

\bibitem{min2023factscore}
Min, S. \textit{et al.} FactScore: Fine-grained Atomic Evaluation of Factual Precision in Long Form Text Generation. \textit{EMNLP}, 2023.

\bibitem{lewis2020retrieval}
Lewis, P. \textit{et al.} Retrieval-Augmented Generation for Knowledge-Intensive NLP Tasks. \textit{NeurIPS}, 2020.

\bibitem{liu2024lost}
Liu, N. F. \textit{et al.} Lost in the Middle: How Language Models Use Long Contexts. \textit{Transactions of the ACL}, 12:157--173, 2024.

\end{thebibliography}

\end{document}
